\documentclass[11pt]{article}

%% This document REQUIRES pdfLaTeX -- see the note above \usepackage{listings}.
%% Compiling with XeLaTeX/LuaLaTeX silently drops Unicode from the Lean code
%% blocks, so fail loudly instead.
\usepackage{iftex}
\RequirePDFTeX

\usepackage{enumitem}
\usepackage{amsthm}
\usepackage{amssymb, amsmath, xcolor}
\usepackage{geometry}
\usepackage[numbers,square,sort&compress]{natbib}
\usepackage{url}
\usepackage[hidelinks]{hyperref}
\usepackage[capitalise,noabbrev]{cleveref}
\usepackage{xspace}
\usepackage{booktabs}
\usepackage{graphicx}
\usepackage{microtype}

\newcommand{\Lean}{\textsc{Lean}~4\xspace}
\newcommand{\Mathlib}{\textsc{Mathlib}\xspace}
\newcommand{\A}{\mathfrak{A}}
\newcommand{\norm}[1]{\left\|#1\right\|}
\newcommand{\abs}[1]{\left|#1\right|}
\newcommand{\comm}[2]{[#1,\,#2]}
\DeclareMathOperator{\tr}{tr}

%% --- Hyperlinked Lean references ---
%% \leanref{File}{line}{name} renders as clickable \texttt{name}
%% linking to GitHub source.  \leansrc{File} links to a whole file.
\newcommand{\repourl}{https://github.com/Jue-Xu/Lean-Trotter/blob/main}
\newcommand{\leanref}[3]{\href{\repourl/LieTrotter/#1\#L#2}{\texttt{#3}}}
\newcommand{\leansrc}[1]{\href{\repourl/LieTrotter/#1}{\texttt{#1}}}

\geometry{margin=1in}

%% Lean identifiers such as \texttt{symmetric\_bch\_septic\_sub\_poly\_axiom} are
%% long, unhyphenatable tokens; give TeX room to reflow around them rather than
%% overflowing the margin.
\setlength{\emergencystretch}{3em}

%% --- Lean code blocks ---
%% COMPILE WITH pdfLaTeX.  (On Overleaf: Menu -> Compiler -> pdfLaTeX.)
%% The Lean source uses Unicode (𝕂 𝔸 ℝ ‖ ≤ • ∑ ⁻¹ τ δ ...), which we render via
%% the `literate` table below.  That mechanism is pdfTeX-only: under XeLaTeX or
%% LuaLaTeX it does not fire, the characters fall through to the monospace font,
%% and every glyph the font lacks is dropped *silently* -- no error, just Lean
%% code blocks with characters missing from the PDF.  The \RequirePDFTeX guard
%% above turns that silent corruption into a loud failure.
%% No -shell-escape is needed (we use `listings`, not `minted`).
\usepackage{listings}
\lstdefinelanguage{Lean4}{
  morekeywords={theorem,lemma,def,noncomputable,section,variable,open,import,
    private,where,by,calc,exact,rw,simp,intro,apply,refine,obtain,have,let,set,
    fun,in,if,then,else,match,with,instance,class,structure,namespace,end,
    include,omit,induction,cases},
  sensitive=true,
  morecomment=[l]{--},
  morecomment=[s]{/-}{-/},
  morestring=[b]",
  literate=
    {𝕂}{{$\mathbb{K}$}}1
    {𝔸}{{$\mathfrak{A}$}}1
    {ℝ}{{$\mathbb{R}$}}1
    {ℂ}{{$\mathbb{C}$}}1
    {ℕ}{{$\mathbb{N}$}}1
    {ℚ}{{$\mathbb{Q}$}}1
    {→}{{$\to$}}1
    {←}{{$\leftarrow$}}1
    {↦}{{$\mapsto$}}1
    {≤}{{$\le$}}1
    {≥}{{$\ge$}}1
    {≠}{{$\ne$}}1
    {⁻¹}{{$^{-1}$}}1
    {•}{{$\bullet$}}1
    {∀}{{$\forall$}}1
    {∃}{{$\exists$}}1
    {∈}{{$\in$}}1
    {∑}{{$\sum$}}1
    {∫}{{$\int$}}1
    {∞}{{$\infty$}}1
    {‖}{{$\|$}}1
    {·}{{$\cdot$}}1
    {α}{{$\alpha$}}1
    {β}{{$\beta$}}1
    {γ}{{$\gamma$}}1
    {δ}{{$\delta$}}1
    {ε}{{$\varepsilon$}}1
    {τ}{{$\tau$}}1
}
\lstset{
  language=Lean4,
  basicstyle=\small\ttfamily,
  keywordstyle=\bfseries\color{blue!70!black},
  commentstyle=\itshape\color{green!50!black},
  stringstyle=\color{red!60!black},
  numbers=left,
  numberstyle=\tiny\color{gray},
  numbersep=6pt,
  frame=lines,
  framesep=2mm,
  breaklines=true,
  tabsize=2,
  showstringspaces=false
}
\lstnewenvironment{leancode}{\lstset{language=Lean4}}{}

%% DO NOT add \DeclareUnicodeCharacter for the astral-plane chars (K/A double-
%% struck, U+1D542 / U+1D538).  Registering them with the UTF-8 decoder collides
%% with the byte-level `literate` machinery above: on the LaTeX 2025-06-01 kernel
%% (which Overleaf runs) every 𝕂/𝔸 inside a listing then dies with
%%   "Invalid UTF-8 byte sequence (\lst@FillFixed@\lst@EC...)".
%% Older kernels (e.g. 2024-11-01) tolerate it, so this does NOT reproduce on a
%% stale local TeX -- it only breaks on Overleaf.  `literate` already renders
%% these characters inside listings, which is the only place they may appear:
%% keep ALL prose ASCII and write $\mathbb{K}$ / $\mathfrak{A}$ there instead.

\title{Formalizing the Lie--Trotter Product Formula and \\
Tightening Commutator-Scaling Error Bounds in Lean~4}
\author{Jue Xu}

\newtheorem{theorem}{Theorem}[section]
\newtheorem{lemma}[theorem]{Lemma}
\newtheorem{proposition}[theorem]{Proposition}
\newtheorem{corollary}[theorem]{Corollary}

\theoremstyle{definition}
\newtheorem{definition}[theorem]{Definition}

\theoremstyle{remark}
\newtheorem*{remark}{Remark}

\begin{document}
\maketitle

\begin{abstract}
We present a machine-checked proof of the Lie--Trotter product
formula in \Lean with \Mathlib---to our knowledge the first
kernel-checked formal proof, with no project-specific axioms, of norm
convergence of the Lie--Trotter sequence for arbitrary elements of a
complete normed algebra---together with higher-order extensions
(symmetric Strang, multi-operator, fourth-order Suzuki $S_4$) and
commutator-scaling error bounds obtained via a Duhamel/FTC
technique that sidesteps ODE uniqueness theory. For the $S_4$
integrator, CAS-assisted Baker--Campbell--Hausdorff expansion yields
four order-$t^5$ bounds, including a replacement of the prefactors of
Childs et al.\ (2021)---rigorous upper bounds not proven tight---by
CAS-certified coefficients that are smaller term by term by factors of
$8.6\times$--$64\times$ (two of the eight vanish identically), the
comparison machine-checked at the leading $t^5$ coefficient; compounding the
single-step bound over $n$ steps yields a machine-checked $O(1/n^4)$
convergence theorem with a commutator-scaled leading coefficient,
completing a verified hierarchy $O(1/n) \to O(1/n^2) \to O(1/n^4)$.
The formalization spans over 14\,000 lines with zero \texttt{sorry}
axioms. At the pinned revision of the companion
\href{https://github.com/Jue-Xu/Lean-BCH}{Lean-BCH} project, every
release headline---including the tight 4th-order Trotter bound and the
local small-$\tau$ refinement at order $\tau^7$---depends only on
\Lean's standard foundational axioms.
\end{abstract}

\tableofcontents
\bigskip

\section{Introduction}
\label{sec:intro}

Product formulas are among the most widely used tools in quantum simulation.
The Lie--Trotter formula
\begin{equation}\label{eq:lie-trotter}
  e^{A+B} = \lim_{n\to\infty} \bigl(e^{A/n}\,e^{B/n}\bigr)^n
\end{equation}
decomposes the exponential of a sum into a product of simpler exponentials,
enabling the simulation of composite quantum systems by sequentially
simulating each term of the Hamiltonian~\cite{trotter1959product, lloydUniversalQuantumSimulators1996}.
Higher-order variants---the symmetric Strang splitting~\cite{strang1968construction}
and Suzuki's recursive hierarchy~\cite{suzukiGeneralTheoryFractal1991}---achieve
faster convergence by exploiting symmetry and commutator cancellation.
Rigorous error bounds for these formulas are critical for resource estimation
in quantum algorithms~\cite{childsTheoryTrotterError2021}.

The quality of these error bounds directly affects the efficiency of quantum
algorithms.  The standard bound scales with the product $\norm{A}\norm{B}$,
which treats $A$ and $B$ as unrelated.
Childs et al.~\cite{childsTheoryTrotterError2021} initiated a program of
\emph{commutator-scaling} bounds, showing that the leading error is controlled
by the commutator norm $\norm{\comm{B}{A}}$ rather than the product
$\norm{A}\norm{B}$.
For spatially local Hamiltonians $H = \sum_j h_j$ typical of many-body physics,
most pairs of terms commute, so $\sum_{j<k} \norm{\comm{h_j}{h_k}}$ can be
polynomially smaller than $\sum_{j<k} \norm{h_j}\norm{h_k}$.
These tighter bounds translate directly into fewer Trotter steps---and hence
fewer quantum gates---for a given simulation accuracy.

In this work, we present a complete formalization of the Lie--Trotter product
formula and its higher-order extensions in \Lean, an interactive theorem prover
with a small trusted kernel~\cite{demoura2021lean4}.
The \Mathlib library provides a large corpus of formalized
mathematics~\cite{mathlib2020}, including the normed algebra exponential,
upon which our development builds.

Proofs of product formula error bounds involve long chains of norm estimates
in non-commutative algebras, where each step interacts with the previous through
submultiplicativity, the triangle inequality, and exponential series manipulation.
A single bookkeeping error---a missing factor, a flipped sign in a
commutator---can silently propagate and invalidate the final bound.
Formal verification forces every intermediate step to be explicit and
machine-checked, making it a natural fit for this domain.

Our formalization proves the following hierarchy of results:
\begin{enumerate}
  \item \textbf{First-order Lie--Trotter} (\cref{sec:assembly}):
    $\norm{(e^{A/n}\,e^{B/n})^n - e^{A+B}} \le C/n$
    with explicit constant $C = 2\norm{A}\norm{B}e^{\norm{A}+\norm{B}} + 1$.

  \item \textbf{Strang splitting} (\cref{sec:strang}):
    $\norm{(e^{A/2n}\,e^{B/n}\,e^{A/2n})^n - e^{A+B}} \le C'/n^2$.

  \item \textbf{Multi-operator extensions}:
    $(e^{A_1/n} \cdots e^{A_m/n})^n \to e^{A_1+\cdots+A_m}$
    and the palindromic (Strang) variant, both with explicit error bounds.

  \item \textbf{Commutator-scaling bounds} (\cref{sec:commutator}):
    \begin{align}
      \norm{e^{tB}\,e^{tA} - e^{t(A+B)}}
        &\le \tfrac{1}{2}\norm{\comm{B}{A}}\,t^2\,e^{t(\norm{A}+3\norm{B})},
        \label{eq:comm-first} \\
      \norm{S_2(t) - e^{tH}}
        &\le \bigl(\tfrac{1}{12}\norm{\comm{B}{\comm{B}{A}}}
          + \tfrac{1}{24}\norm{\comm{A}{\comm{A}{B}}}\bigr)\,t^3,
        \label{eq:comm-strang}
    \end{align}
    matching the commutator-scaling program of Childs et al.~\cite{childsTheoryTrotterError2021}.

  \item \textbf{Tighter Strang bound via norm-of-difference}: a strictly
    sharper second-order bound
    $\norm{S_2(t)-e^{tH}} \le \tfrac{\norm{D}}{6}t^3 + \tfrac{T}{4}t^4$
    where $D = \comm{B}{\comm{B}{A/2}} - \comm{A/2}{\comm{A/2}{B}}$ is an
    effective double commutator capturing signed cancellations that the
    triangle-inequality form of \eqref{eq:comm-strang} cannot see.

  \item \textbf{Fourth-order Suzuki $S_4$ with CAS-assisted BCH bounds}
    (\cref{sec:suzuki4}):
    four formalized order-$t^5$ bounds for
    $\norm{S_4(t) - e^{t(A+B)}}$, culminating in a provably
    strictly tighter (9$\times$--64$\times$ per coefficient)
    replacement of Childs's~2021 heuristic prefactors, together with an
    existential-$\delta$ finite-$t$ bound obtained by extending the CAS
    computation to order~$\tau^7$.

  \item \textbf{Fourth-order convergence of $S_4$} (\cref{sec:s4-convergence}):
    compounding the single-step $O(t^5)$ bound over $n$ steps gives
    $\norm{(S_4(t/n))^n - e^{t(A+B)}} \le C/n^4$ for $n$ beyond an explicit
    threshold, and hence $(S_4(t/n))^n \to e^{t(A+B)}$. Together with items~1
    and~2 this completes a verified convergence hierarchy
    $O(1/n) \to O(1/n^2) \to O(1/n^4)$.
\end{enumerate}

The entire development comprises 32 Lean source files totaling over 12\,100 lines,
with zero \texttt{sorry} axioms and zero own theorem-level BCH-interface
axioms on the Lean-Trotter side. The three former BCH-interface axioms
(\texttt{bch\_w4Deriv\_quintic\_level2},
\texttt{bch\_w4Deriv\_level3\_tight},
\texttt{bch\_uniform\_integrated})
are now theorems composing Lean-BCH bridge corollaries with
exp-Lipschitz lifts. As of Lean-BCH revision \texttt{d455ff0}
(2026-05-19), the upstream B1.c quintic axiom has been discharged, so
all $\tau^5$ Suzuki~$S_4$ headlines---together with the standard
Lie--Trotter, Strang, and commutator-scaling theorems---depend only
on the standard foundational axioms \texttt{[propext, Classical.choice,
Quot.sound]}. In particular, axiom inspection on
\texttt{norm\_suzuki4\_childs\_form\_via\_level3} (Level~1 Childs
reproduction), \texttt{norm\_suzuki4\_level3\_bch} (Level~3 tight
$\gamma_i$ prefactors), \texttt{norm\_suzuki4\_level2\_bch} (Level~2
unit-coefficient $\tau^5$ bound), the abstract $S_4$ $O(t^5)$ identity
(\texttt{bch\_iteratedDeriv\_s4Func\_order4}), and \texttt{lie\_trotter}
returns the three foundational axioms only. Two transitive private
axioms remain inside the companion
\href{https://github.com/Jue-Xu/Lean-BCH}{Lean-BCH} project---both
septic stepping stones
(\texttt{symmetric\_bch\_septic\_sub\_poly\_axiom},
\texttt{norm\_septic\_match\_residual\_le\_axiom}). They gate only
the optional Level~4 finite-$t$ uniform refinement with an explicit
$\tau^7$ correction, not the core tight 4th-order bound. All other
Lean-BCH content is fully sorry-free; the three-factor symmetric BCH
cubic identities, the two-factor palindromic quintic remainder, and
the five-factor palindromic quintic identification are theorems
specialized to $\mathbb{K} = \mathbb{R}$ and imported directly.
Our \Lean formalization is available at
\href{https://github.com/Jue-Xu/Lean-Trotter}{\texttt{github.com/Jue-Xu/Lean-Trotter}}.
Our main contributions are:
\begin{itemize}
  \item The \emph{first formalization} of the Lie--Trotter product formula in any
    proof assistant, including its higher-order and multi-operator extensions.

  \item A \emph{Duhamel/FTC-based proof technique} for commutator-scaling bounds
    that avoids ODE uniqueness theory (Gronwall's inequality, Picard iteration),
    requiring only the fundamental theorem of calculus for Bochner integrals.

  \item Machine-checked \emph{commutator-scaling error bounds} at first and
    second order, with double-commutator coefficients matching the
    theoretical literature---and a new tighter-Strang refinement
    via norm-of-difference (\cref{apd:tighter}).

  \item \emph{CAS-assisted BCH bounds for Suzuki $S_4$}: four machine-checked
    bounds---including a rigorously tighter replacement for the
    Childs (2021) prefactors and an existential-$\delta$ finite-$t$ bound
    via an order-$\tau^7$ correction---with all CAS output verified by Lean
    numerical sanity checks.

  \item Identification and filling of several \emph{gaps in \Mathlib}, notably
    the inequality $\norm{e^a} \le e^{\norm{a}}$ for general Banach algebras.
\end{itemize}

To our knowledge, no prior formalization of the Lie--Trotter product formula
exists in any proof assistant.
\Mathlib contains the identity $e^{a+b} = e^a\,e^b$ for commuting elements
(\texttt{exp\_add\_of\_commute}), but the non-commuting convergence
theorem---which is a limit statement, not an algebraic identity---is new.

The remainder of this paper is organized as follows.
\Cref{sec:background} reviews the mathematical background on product formulas
and the Lean/\Mathlib ecosystem.
\Cref{sec:proof} presents the proof architecture, including the modular
decomposition and the Duhamel technique for commutator scaling.
\Cref{sec:formalization} discusses the formalization in \Lean, with
representative code excerpts and design decisions.
\Cref{sec:discussion} reflects on the development process, including
AI-assisted formalization and lessons learned.
\Cref{sec:conclusion} concludes with future directions.

\section{Mathematical Background}
\label{sec:background}

\subsection{The Lie--Trotter product formula}
\label{sec:bg-trotter}

Let $\A$ be a unital Banach algebra over $\mathbb{K} \in \{\mathbb{R}, \mathbb{C}\}$
with norm $\norm{\cdot}$ satisfying $\norm{1} = 1$ and $\norm{ab} \le \norm{a}\norm{b}$.
The exponential
\[
  e^a = \sum_{k=0}^{\infty} \frac{a^k}{k!}
\]
converges absolutely for every $a \in \A$.

Trotter~\cite{trotter1959product} proved that for bounded operators $A, B$,
\[
  e^{A+B} = \lim_{n\to\infty} \bigl(e^{A/n}\,e^{B/n}\bigr)^n.
\]
The standard proof proceeds in two steps:
\begin{enumerate}
  \item \emph{Telescoping}: writing $M = \max\{\norm{X},\norm{Y}\}$,
    the identity
    $X^n - Y^n = \sum_{k=0}^{n-1} X^k(X-Y)Y^{n-1-k}$ gives
    $\norm{X^n - Y^n} \le n\,\norm{X-Y}\,M^{n-1}$.

  \item \emph{Step error}: the single-step approximation satisfies
    $\norm{e^a\,e^b - e^{a+b}} = O(\norm{a}\,\norm{b})$.
\end{enumerate}
Setting $X = e^{A/n}\,e^{B/n}$ and $Y = e^{(A+B)/n}$, the $O(1/n^2)$ step
error is amplified by the factor $n$ from telescoping, yielding $O(1/n)$
overall convergence.

\subsection{Higher-order splitting methods}
\label{sec:bg-splitting}

The symmetric \emph{Strang splitting}~\cite{strang1968construction} uses
\[
  S_2(t) = e^{tA/2}\,e^{tB}\,e^{tA/2}
\]
to achieve $O(t^3)$ local error and $O(1/n^2)$ global convergence.
Suzuki~\cite{suzukiGeneralTheoryFractal1991} showed that higher-order integrators can be
constructed recursively: the fourth-order formula
\[
  S_4(t) = S_2(pt)^2\,S_2\bigl((1-4p)t\bigr)\,S_2(pt)^2
  \qquad \text{with } p = \frac{1}{4 - 4^{1/3}}
\]
achieves $O(t^5)$ local error.

Both formulas generalize to multiple operators $A_1, \ldots, A_m$.
The first-order product $(e^{A_1/n} \cdots e^{A_m/n})^n$ converges at rate
$O(1/n)$, while the palindromic product
\[
  S_2^{(m)}(t) = e^{tA_1/2}\,e^{tA_2/2}\cdots e^{tA_m}\cdots e^{tA_2/2}\,e^{tA_1/2}
\]
achieves $O(1/n^2)$.

\subsection{Commutator-scaling error bounds}
\label{sec:bg-commutator}

The standard step error $\norm{e^a\,e^b - e^{a+b}} \le 2\norm{a}\norm{b}e^{\norm{a}+\norm{b}}$
treats $A$ and $B$ as unrelated. When $A$ and $B$ nearly commute---as is
typical for spatially local Hamiltonians---the error is much smaller.

Childs et al.~\cite{childsTheoryTrotterError2021} showed that the leading error term is
controlled by the commutator $\comm{B}{A} = BA - AB$ rather than the product
$\norm{A}\norm{B}$. At first order,
\begin{equation}\label{eq:comm-scaling-bg}
  \norm{e^{tB}\,e^{tA} - e^{t(A+B)}}
    \le \frac{\norm{\comm{B}{A}}}{2}\,t^2\,e^{t(\norm{A}+3\norm{B})}.
\end{equation}
For the Strang splitting, the leading commutator $\comm{B}{A}$ cancels by
symmetry, and the error is controlled by double commutators:
\begin{equation}\label{eq:strang-comm-bg}
  \norm{S_2(t) - e^{tH}}
    \le \Bigl(\frac{\norm{\comm{B}{\comm{B}{A}}}}{12}
      + \frac{\norm{\comm{A}{\comm{A}{B}}}}{24}\Bigr)\,t^3.
\end{equation}

These commutator-scaling bounds are physically significant: for a lattice
Hamiltonian $H = \sum_j h_j$ where each $h_j$ acts on a few neighboring sites,
$\norm{\comm{h_j}{h_k}} = 0$ whenever the supports of $h_j$ and $h_k$ are
disjoint, so the commutator-based bound can be polynomially tighter than the
naive product bound~\cite{childsTheoryTrotterError2021}.

\subsection{Lean~4 and Mathlib}
\label{sec:bg-lean}

\Lean is an interactive theorem prover and programming language based on
dependent type theory~\cite{demoura2021lean4}. Proofs and programs share the
same language, and correctness is guaranteed by a small trusted kernel that
independently type-checks every proof term.

The \Mathlib library~\cite{mathlib2020} provides formalized mathematics
including abstract algebra, analysis, and topology. For our development,
the key components are:
\begin{itemize}
  \item \texttt{NormedRing}, \texttt{NormedAlgebra}: typeclasses for
    normed rings and algebras with submultiplicative norms.
  \item \texttt{NormedSpace.exp}: the exponential function
    $a \mapsto \sum_{k} a^k/k!$ in a complete normed algebra.
  \item \texttt{intervalIntegral}: the Bochner integral over $[a,b]$,
    with the fundamental theorem of calculus.
  \item \texttt{HasDerivAt}: the Fr\'{e}chet derivative, with product rules
    and chain rules for normed algebra-valued functions.
\end{itemize}

\paragraph{Related formalization work.}
Formal verification has recently gained traction in mathematical physics and
quantum information.
Zhao and Yu~\cite{zhao2025chsh} formalized robust CHSH rigidity in \Lean,
discovering a gap in the original proof of McKague~(2012) during the
formalization process.
Meiburg et al.~\cite{meiburg2025steinlemma} formalized a generalized
quantum Stein's lemma, and
Tooby-Smith~\cite{toobysmith2026higgs} formalized the stability of the
two-Higgs-doublet potential, also identifying a literature error.
These projects share the pattern of using formal verification not merely to
certify known results but to expose subtle gaps in published proofs.
Our work extends this program to quantum simulation, providing the first
machine-checked proof of any product formula for operator exponentials.

\section{Proof Architecture}
\label{sec:proof}

Our formalization is organized into five independent tracks that compose to yield
the main theorem and its extensions. \Cref{fig:dag} shows the dependency
structure. We present the mathematical arguments here; the Lean formalization
is discussed in \cref{sec:formalization}.

\begin{figure}[ht]
\centering
\small
\[
\underbrace{
  \begin{array}{c}
  \texttt{Telescoping} \\[2pt]
  X^n - Y^n = \sum X^k(X{-}Y)Y^{n{-}1{-}k}
  \end{array}
}_{\text{Track 1}}
\quad
\underbrace{
  \begin{array}{c}
  \texttt{ExpBounds} \\[2pt]
  \|e^a\| \le e^{\|a\|}
  \end{array}
}_{\text{Track 2}}
\quad
\underbrace{
  \begin{array}{c}
  \texttt{ExpDivPow} \\[2pt]
  (e^{a/n})^n = e^a
  \end{array}
}_{\text{Track 4}}
\]
\[
\downarrow \qquad\qquad\qquad\qquad \downarrow
\]
\[
\underbrace{
  \begin{array}{c}
  \texttt{StepError} \\[2pt]
  \|e^a e^b - e^{a+b}\| \le 2\|a\|\|b\|\,e^{\|a\|+\|b\|}
  \end{array}
}_{\text{Track 3 (depends on Track 2)}}
\]
\[
\downarrow
\]
\[
\boxed{
  \begin{array}{c}
  \texttt{Assembly} \\[2pt]
  (e^{A/n}\,e^{B/n})^n \to e^{A+B} \quad [O(1/n)]
  \end{array}
}
\]
\[
\swarrow \qquad\qquad\qquad \searrow
\]
\[
\underbrace{
  \begin{array}{c}
  \texttt{StrangSplitting} \\[2pt]
  O(1/n^2)
  \end{array}
}_{\text{Extensions}}
\quad
\underbrace{
  \begin{array}{c}
  \texttt{Multi\{Operator,Strang\}} \\[2pt]
  m\text{-operator}
  \end{array}
}_{\text{Extensions}}
\quad
\underbrace{
  \begin{array}{c}
  \texttt{CommutatorScaling} \\[2pt]
  \|[B,A]\|\text{-scaling}
  \end{array}
}_{\text{Track 5 (independent)}}
\]
\caption{Dependency structure of the formalization.
  Tracks~1--4 compose into the main theorem; extensions branch from the
  assembly or develop independently.}
\label{fig:dag}
\end{figure}

\subsection{Modular decomposition}
\label{sec:modular}

The proof of the Lie--Trotter formula decomposes into four independent
components:
\begin{description}[leftmargin=1.5em]
  \item[Track~1 (Algebraic):] The telescoping identity and its norm bound.
  \item[Track~2 (Analytic):] Exponential series remainder estimates.
  \item[Track~3 (Step error):] The single-step approximation bound
    $\norm{e^a\,e^b - e^{a+b}} \le 2\norm{a}\norm{b}\,e^{\norm{a}+\norm{b}}$.
  \item[Track~4 (Connecting):] The identity $(e^{a/n})^n = e^a$.
\end{description}
Tracks 1--4 are combined in an assembly step to yield the convergence rate,
from which the limit theorem follows by a standard $\varepsilon$-$N$ argument.
The commutator-scaling bounds (Track~5) are an independent extension that
refines the step error.

\subsection{Telescoping and norm bounds}
\label{sec:telescoping}

The algebraic identity
\begin{equation}\label{eq:telescoping}
  X^n - Y^n = \sum_{k=0}^{n-1} X^k\,(X-Y)\,Y^{n-1-k}
\end{equation}
is proved by induction on $n$. The key step is factoring $Y$ out of the inner
sum to match the inductive hypothesis: writing $Y^{m-k} = Y^{m-1-k} \cdot Y$
converts $\sum_{k<m} X^k(X-Y)Y^{m-k}$ into $(\sum_{k<m} X^k(X-Y)Y^{m-1-k}) \cdot Y$.

The norm bound follows by the triangle inequality applied to each summand:
\begin{equation}\label{eq:norm-telescoping}
  \norm{X^n - Y^n} \le n\,\norm{X-Y}\,M^{n-1},
  \qquad M = \max\{\norm{X}, \norm{Y}\}.
\end{equation}

\subsection{Exponential series estimates}
\label{sec:exp-bounds}

We establish four bounds on the exponential in a complete normed algebra:
\begin{align}
  \norm{e^a} &\le e^{\norm{a}},
    \label{eq:norm-exp} \\
  \norm{e^a - 1} &\le e^{\norm{a}} - 1,
    \label{eq:norm-exp-sub-one} \\
  e^x - 1 - x &\le \tfrac{x^2}{2}\,e^x
    \quad (x \ge 0),
    \label{eq:exp-quadratic} \\
  \norm{e^a - 1 - a} &\le \tfrac{\norm{a}^2}{2}\,e^{\norm{a}}.
    \label{eq:norm-exp-quadratic}
\end{align}

The proofs use the power series representation $e^a = \sum_{n=0}^{\infty} a^n/n!$.
Bounds \eqref{eq:norm-exp}--\eqref{eq:norm-exp-sub-one} follow from
$\norm{\sum a^n/n!} \le \sum \norm{a}^n/n! = e^{\norm{a}}$ with appropriate
index shifts. Bound \eqref{eq:exp-quadratic} uses the termwise estimate
$x^{n+2}/(n+2)! \le (x^2/2) \cdot x^n/n!$ via the factorial inequality
$2 \cdot n! \le (n+2)!$. Bound \eqref{eq:norm-exp-quadratic} combines
\eqref{eq:exp-quadratic} with the same tsum comparison technique as
\eqref{eq:norm-exp}.

\begin{remark}
  Bound \eqref{eq:norm-exp} was not available in \Mathlib for general
  Banach algebras; \Mathlib only had the specialized result
  \texttt{Complex.norm\_exp\_le\_exp\_norm} for $\mathbb{C}$.
  Our proof from the power series fills this gap.
\end{remark}

\subsection{Step error via algebraic factorization}
\label{sec:step-error}

The single-step bound is the most delicate component. Rather than expanding
$e^a\,e^b$ and $e^{a+b}$ to second order and bounding cross-terms directly
(which is tractable on paper but unwieldy in a proof assistant due to
non-commutative bookkeeping), we use an algebraic factorization:
\begin{equation}\label{eq:factorization}
  e^a\,e^b - e^{a+b}
    = (e^a - 1)(e^b - 1) - \bigl(e^{a+b} - e^a - e^b + 1\bigr).
\end{equation}
The first term is bounded by $(e^{\norm{a}}-1)(e^{\norm{b}}-1)$ via
\eqref{eq:norm-exp-sub-one}. For the second (cross) term, we prove by
induction on the power $m$ that
\begin{equation}\label{eq:cross-term}
  \norm{(a+b)^m - a^m - b^m}
    \le (\norm{a}+\norm{b})^m - \norm{a}^m - \norm{b}^m
    \qquad (m \ge 1),
\end{equation}
using the non-commutative recurrence
$(a+b)^{m+1} - a^{m+1} - b^{m+1} = (a+b)\bigl((a+b)^m - a^m - b^m\bigr) + ab^m + ba^m$.
Summing over the exponential series, the cross term is also bounded by
$(e^{\norm{a}}-1)(e^{\norm{b}}-1)$, so by the triangle inequality:
\begin{equation}\label{eq:step-bound}
  \norm{e^a\,e^b - e^{a+b}}
    \le 2(e^{\norm{a}}-1)(e^{\norm{b}}-1)
    \le 2\norm{a}\norm{b}\,e^{\norm{a}+\norm{b}},
\end{equation}
where the last step uses $(e^s-1)(e^t-1) \le st\,e^{s+t}$ for $s,t \ge 0$.

\begin{remark}
  The factorization \eqref{eq:factorization} was discovered during
  formalization: the standard second-order expansion approach required
  tracking non-commutative multinomial cross-terms, which was intractable
  in Lean. The algebraic factorization reduces the problem to two
  independent bounds, each handled by series comparison.
\end{remark}

\subsection{Assembly and convergence}
\label{sec:assembly}

Setting $P = e^{A/n}\,e^{B/n}$ and $Q = e^{(A+B)/n}$:
\begin{enumerate}
  \item $Q^n = e^{A+B}$ by the identity $(e^{a/n})^n = e^a$
    (\leanref{ExpDivPow.lean}{31}{exp\_div\_pow}).
  \item $\norm{P-Q} \le 2\norm{A}\norm{B}/n^2 \cdot e^{(\norm{A}+\norm{B})/n}$
    by the step error with $a = A/n$, $b = B/n$.
  \item $\max\{\norm{P},\norm{Q}\} \le e^{(\norm{A}+\norm{B})/n}$ by
    \eqref{eq:norm-exp}.
  \item Telescoping gives
    $\norm{P^n - Q^n} \le n \cdot \norm{P-Q} \cdot M^{n-1}$.
  \item The factors $e^{(\norm{A}+\norm{B})/n}$ combine:
    $e^{s/n} \cdot (e^{s/n})^{n-1} = (e^{s/n})^n = e^s$, yielding
    \[
      \norm{(e^{A/n}\,e^{B/n})^n - e^{A+B}}
        \le \frac{2\norm{A}\norm{B}\,e^{\norm{A}+\norm{B}}}{n}.
    \]
\end{enumerate}
The main theorem $\lim_{n\to\infty} (e^{A/n}\,e^{B/n})^n = e^{A+B}$ follows
by a standard $\varepsilon$-$N$ argument using the Archimedean property.

\subsection{Duhamel approach to commutator scaling}
\label{sec:commutator}

The commutator-scaling bound \eqref{eq:comm-first} replaces the product
$\norm{A}\norm{B}$ with $\norm{\comm{B}{A}}$. Our proof uses a
variation-of-parameters (Duhamel) technique based on the fundamental theorem
of calculus, avoiding ODE uniqueness theory entirely.

\paragraph{Integral error representation.}
Define the ``conjugated Trotter product''
\begin{equation}\label{eq:w-def}
  w(\tau) = e^{-\tau(A+B)}\,e^{\tau B}\,e^{\tau A}.
\end{equation}
By the product rule for Fr\'{e}chet derivatives in a normed algebra,
\begin{equation}\label{eq:w-deriv}
  w'(\tau) = e^{-\tau(A+B)}\,\bigl(e^{\tau B}\,A - A\,e^{\tau B}\bigr)\,e^{\tau A}
  = e^{-\tau(A+B)}\,\comm{e^{\tau B}}{A}\,e^{\tau A}.
\end{equation}
The key cancellation is that the ``free'' terms involving $A+B$, $B$, and $A$
sum to zero: $-(A+B) + B + A = 0$. The FTC-2 gives
\begin{equation}\label{eq:integral-error}
  e^{tB}\,e^{tA} - e^{t(A+B)}
    = \int_0^t e^{(t-\tau)(A+B)}\,\comm{e^{\tau B}}{A}\,e^{\tau A}\,d\tau,
\end{equation}
after left-multiplying by $e^{t(A+B)}$ (which commutes with itself).

\paragraph{Extracting the commutator.}
The integrand contains $\comm{e^{\tau B}}{A} = e^{\tau B}A - Ae^{\tau B}$,
which involves the \emph{exponential} of $B$. To extract the bare commutator
$\comm{B}{A}$, we apply FTC-2 a second time to the conjugation map
$s \mapsto e^{sB}\,A\,e^{-sB}$:
\begin{equation}\label{eq:conj-ftc}
  e^{\tau B}\,A\,e^{-\tau B} - A
    = \int_0^{\tau} e^{sB}\,\comm{B}{A}\,e^{-sB}\,ds.
\end{equation}
This gives $\norm{e^{\tau B}\,A\,e^{-\tau B} - A}
  \le \norm{\comm{B}{A}}\,\abs{\tau}\,e^{2\abs{\tau}\norm{B}}$
by bounding the integrand and using
$\norm{\int_0^\tau f} \le C\abs{\tau}$.
Right-multiplying by $e^{\tau B}$ yields
$\norm{\comm{e^{\tau B}}{A}} \le \norm{\comm{B}{A}}\,\abs{\tau}\,e^{3\abs{\tau}\norm{B}}$.

\paragraph{Final bound.}
Substituting into \eqref{eq:integral-error} and bounding the integral:
\[
  \norm{e^{tB}\,e^{tA} - e^{t(A+B)}}
    \le \int_0^t e^{(t-\tau)(\norm{A}+\norm{B})} \cdot
      \norm{\comm{B}{A}}\,\tau\,e^{3\tau\norm{B}} \cdot e^{\tau\norm{A}}\,d\tau.
\]
Simplifying the exponential factors and evaluating $\int_0^t \tau\,d\tau = t^2/2$
gives the claimed bound \eqref{eq:comm-first}.

\begin{remark}
  The Duhamel approach requires only the product rule for Fr\'{e}chet
  derivatives, the FTC-2 for Bochner integrals, and basic norm estimates---all
  available in \Mathlib. The alternative approach via Gronwall's inequality
  would require ODE existence/uniqueness theory, adding approximately 40 lines
  of additional Lean infrastructure.
\end{remark}

\subsection{Second-order Strang commutator scaling}
\label{sec:strang}

For the Strang splitting $S_2(t) = e^{tA/2}\,e^{tB}\,e^{tA/2}$, we prove
\begin{equation}\label{eq:strang-bound}
  \norm{S_2(t) - e^{tH}}
    \le \frac{\norm{\comm{B}{\comm{B}{A}}}}{12}\,t^3
      + \frac{\norm{\comm{A}{\comm{A}{B}}}}{24}\,t^3,
\end{equation}
where $H = A + B$ and we additionally require that $A$ and $B$ are
\emph{skew-adjoint} ($a^* = -a$), so that $\norm{e^{ta}} = 1$.

The proof applies the same Duhamel strategy to
$w(\tau) = e^{-\tau H}\,S_2(\tau)$, but now requires:
\begin{enumerate}
  \item A \emph{4-factor product rule} for $S_2(\tau) = e^{\tau A/2}\,e^{\tau B}\,e^{\tau A/2}$,
    yielding $w'(\tau) = e^{-\tau H}\,\mathcal{T}_2(\tau)\,S_2(\tau)$, where $\mathcal{T}_2$
    involves the commutators $\comm{A/2}{e^{\tau B}}$ and $\comm{e^{\tau B}}{A/2}$.

  \item \emph{First-order cancellation}: the leading commutator terms cancel
    exactly: $\comm{A/2}{B} + \comm{B}{A/2} = 0$.

  \item A \emph{double FTC} on $s \mapsto e^{sB}\,A\,e^{-sB}$, expanding to
    second order to extract $\comm{B}{\comm{B}{A}}$.

  \item The \emph{``subtract-constant-at-$\tau$'' trick}: to bound
    $\int_0^\tau F(s)\,ds - \tau\,F(\tau)$, rewrite as
    $\int_0^\tau (F(s) - F(\tau))\,ds$ and bound $\norm{F(s) - F(\tau)}
    \le C_A(\tau - s)$.
    This avoids Fubini's theorem or integration by parts.
\end{enumerate}

\paragraph{The 4-factor product rule.}
Write $A' = A/2$ and $H = A + B$. The conjugated Strang product
\[
  w(\tau) = e^{-\tau H}\,e^{\tau A'}\,e^{\tau B}\,e^{\tau A'}
\]
has four exponential factors, each contributing a derivative term via the
Leibniz rule. Using $d/d\tau[e^{\tau X}] = X\,e^{\tau X}$ and collecting:
\begin{equation}\label{eq:4factor}
  w'(\tau) = e^{-\tau H}\,\mathcal{T}_2(\tau)\,S_2(\tau),
\end{equation}
where the \emph{Strang residual} is
\[
  \mathcal{T}_2(\tau)
    = \bigl(e^{\tau A'}\,B\,e^{-\tau A'} - B\bigr)
      + e^{\tau A'}\bigl(e^{\tau B}\,A'\,e^{-\tau B} - A'\bigr)e^{-\tau A'}.
\]
The key algebraic step is verifying that the ``free'' derivative terms
$-H + A' + B + A'$ sum to zero (since $A' + A' = A$ and $H = A + B$).
In Lean, the sum of the four derivative contributions is factored as
$-(e^{-\tau H}) \cdot (H - A' - A' - B) \cdot S_2(\tau)$, where
$H - A' - A' - B = 0$ is discharged by \texttt{abel}. The free-ring
factoring is handled by \texttt{noncomm\_ring}.

\paragraph{Coefficient conversion.}
Since $A' = A/2$, the double commutators in the final bound involve
half-scaled operators:
$\comm{B}{\comm{B}{A'}} = \frac{1}{2}\comm{B}{\comm{B}{A}}$ and
$\comm{A'}{\comm{A'}{B}} = \frac{1}{4}\comm{A}{\comm{A}{B}}$.
The scalar factors are extracted through $\comm{cX}{Y} = c\,\comm{X}{Y}$
(valid for central scalars) using \texttt{Algebra.smul\_mul\_assoc} and
\texttt{Algebra.mul\_smul\_comm}, then \texttt{norm\_smul}.
Combined with the integration constants ($1/6$ from the cubic integral
and $1/2$ from the double FTC), this produces the final coefficients
$1/12$ and $1/24$ in \eqref{eq:strang-bound}.

\section{Formalization in Lean~4}
\label{sec:formalization}

\subsection{Project structure and statistics}
\label{sec:project-structure}

The formalization is organized into 41~Lean source files, grouped by track
in \cref{tab:files}. Each file is self-contained with explicit imports,
enabling parallel development and independent verification.

\begin{table}[ht]
\centering
\footnotesize
\setlength{\tabcolsep}{4pt}
\begin{tabular}{@{}lrl@{}}
\toprule
\textbf{File} & \textbf{Lines} & \textbf{Content} \\
\midrule
\multicolumn{3}{@{}l}{\textit{Lie--Trotter core (Tracks 1--4, \cref{sec:proof})}} \\
\leansrc{Telescoping.lean}       & 106  & Algebraic telescoping + norm bound \\
\leansrc{ExpBounds.lean}         & 212  & Exponential series remainder estimates \\
\leansrc{StepError.lean}         & 674  & Quadratic step error (algebraic factorization) \\
\leansrc{ExpDivPow.lean}         & 54   & $(e^{a/n})^n = e^a$ \\
\leansrc{Assembly.lean}          & 146  & $O(1/n)$ convergence + main theorem \\
\midrule
\multicolumn{3}{@{}l}{\textit{Strang, multi-operator, commutator scaling (Track~5)}} \\
\leansrc{StrangSplitting.lean}   & 492  & Strang splitting, $O(1/n^2)$ \\
\leansrc{MultiOperator.lean}     & 288  & Multi-operator first-order \\
\leansrc{MultiStrang.lean}       & 428  & Multi-operator Strang \\
\leansrc{CommutatorScaling.lean} & 379  & Duhamel commutator-scaling \\
\leansrc{MultiCommutatorScaling.lean}       & 128  & Multi-operator comm.\ scaling \\
\leansrc{StrangCommutatorScaling.lean}      & 733  & Strang double-commutator scaling \\
\leansrc{MultiStrangCommutatorScaling.lean} & 174  & Multi-op.\ Strang comm.\ scaling \\
\leansrc{HigherCommutator.lean}            & 243  & Triple-FTC: extracts $[B,[B,[B,A]]]$ \\
\leansrc{StrangCommutatorScalingTight.lean} & 603 & Tighter Strang via norm-of-difference $D$; $\norm{D}$ domination \\
\leansrc{StrangTotalErrorCommScaling.lean}  & 184 & Commutator-scaled Strang \emph{total} error \\
\midrule
\multicolumn{3}{@{}l}{\textit{Suzuki $S_4$ track (\cref{sec:suzuki4})}} \\
\leansrc{Suzuki4.lean}                    & 466  & $S_4$ definition; generic $O(1/n^2)$ rate \\
\leansrc{Suzuki4FullDuhamel.lean}         & 445  & $S_4$ $O(t^3)$ via 5-$S_2$ telescoping \\
\leansrc{Suzuki4CommutatorScaling.lean}   & 90   & \texttt{suzuki4Exp} definition \\
\leansrc{Suzuki4Deriv.lean}               & 96   & Derivative scaffolding \\
\leansrc{Suzuki4HasDerivAt.lean}          & 136  & Module~1: \texttt{HasDerivAt} for 12-factor $w_4$ \\
\leansrc{Suzuki4Module2.lean}             & 167  & Module~2: FTC-2 bridge \\
\leansrc{Suzuki4Module3.lean}             & 184  & Module~3: FTC-2 reduction \\
\leansrc{Suzuki4Module4.lean}             & 175  & Module~4a: continuity of $w_4'$ \\
\leansrc{Suzuki4DerivExplicit.lean}       & 979  & Module~4b: explicit derivative; Phase 1--3 \\
\leansrc{Suzuki4StrangBlocks.lean}        & 120  & $S_4$ = 5 Strang blocks + cubic cancellation \\
\leansrc{Suzuki4MultinomialExpand.lean}   & 787  & Multinomial formulas; h2 unconditional, h3 under Suzuki cubic \\
\leansrc{Suzuki4Phase5.lean}              & 782  & Taylor reduction + Leibniz bridges + CAPSTONE \\
\leansrc{Suzuki4OrderFive.lean}           & 440  & $O(t^5)$ abstract-form target \\
\leansrc{Suzuki4ChildsForm.lean}          & 216  & Childs Prop pf4\_bound\_2term framework \\
\leansrc{Suzuki4BchBound.lean}            & 495  & SLICE~1: single-step $O(\abs{\tau}^5)$ bound \\
\leansrc{TaylorMatch.lean}                & 331  & SLICE~2: generic Taylor-match-from-norm \\
\leansrc{Suzuki4ViaBCH.lean}              & 1\,436 & SLICE~3 wiring + L1/L2/L3/L4 BCH-derived bounds \\
\leansrc{Suzuki4Convergence.lean}         & 398  & Total error $O(1/n^4)$ + convergence \\
\leansrc{Suzuki4TightConvergence.lean}    & 169  & Commutator-scaled $O(1/n^4)$ \\
\leansrc{Suzuki4UnitaryTotalError.lean}   & 274  & Skew-adjoint total error, no growth factor \\
\leansrc{Suzuki4Commute.lean}             & 228  & Commuting case: $S_4$ exact, error $\equiv 0$ \\
\leansrc{Suzuki4GapClosers.lean}          & 205  & General-$t$ step; imaginary-time corollaries \\
\leansrc{TrotterStepCount.lean}           & 392  & $\varepsilon$-form step counts; explicit per-$n$ bounds \\
\leansrc{MatrixCorollaries.lean}          & 157  & Matrix specializations (spectral norm) \\
\leansrc{PrefactorStrict.lean}            & 111  & Strict $\gamma_i < \alpha_i$; gap $\gamma_i \le \alpha_i/8$ \\
\midrule
\multicolumn{3}{@{}l}{\textit{Root}} \\
\href{\repourl/LieTrotter.lean}{\texttt{LieTrotter.lean}} & 73 & Import aggregator \\
\midrule
\textbf{Total} & \textbf{14\,196} & \textbf{0 \texttt{sorry} axioms,
0 own theorem-level BCH-interface axioms} \\
\bottomrule
\end{tabular}
\caption{Source files in the formalization, grouped by track.
\texttt{Suzuki4.lean} carries the generic $O(1/n^2)$ rate, which holds for
every $n \ge 1$ and needs no cubic condition;
\texttt{Suzuki4Convergence.lean} upgrades the \emph{same} object to $O(1/n^4)$
under the Suzuki cubic condition (\cref{sec:s4-convergence}), and
\texttt{Suzuki4TightConvergence.lean} does so with a commutator-scaled leading
coefficient.}
\label{tab:files}
\end{table}

The development builds against Lean~4.29.0-rc8 and current \Mathlib.
The companion
\href{https://github.com/Jue-Xu/Lean-BCH}{Lean-BCH} project is imported
as a git dependency, supplying the symmetric three-factor BCH identities,
the two-factor palindromic quintic remainder, and the five-factor
palindromic quintic identification used in the Suzuki~$S_4$ track.
Lean-BCH is sorry-free as of revision~\texttt{d455ff0} (2026-05-19),
which closed the previous B1.c (\texttt{symmetric\_bch\_quintic\_sub\_poly\_axiom})
dependency; its single-step $O(\abs{\tau}^5)$ payoff lemma anchors
our SLICE-1 bound in \leansrc{Suzuki4BchBound.lean}. Two private
septic axioms remain in Lean-BCH and gate only the optional Level~4
local small-$\tau$ ($\tau^7$) refinement (see \cref{sec:caveats}).
Build time for the full project is under 2~minutes on a standard
laptop after \Mathlib oleans are cached (\texttt{lake exe cache get}).

\subsection{Type-level setup}
\label{sec:type-setup}

Our theorems are stated for a general complete normed algebra $\A$ over a
field $\mathbb{K} \in \{\mathbb{R}, \mathbb{C}\}$:

\begin{leancode}
variable {𝕂 : Type*} [RCLike 𝕂]
variable {𝔸 : Type*} [NormedRing 𝔸] [NormedAlgebra 𝕂 𝔸]
  [NormOneClass 𝔸] [CompleteSpace 𝔸]
\end{leancode}

The typeclass \texttt{NormOneClass} requires $\norm{1} = 1$, which is needed
for \texttt{norm\_pow\_le} in current \Mathlib. The main theorem is
(\leanref{Assembly.lean}{121}{lie\_trotter}):

\begin{leancode}
theorem lie_trotter (A B : 𝔸) :
    Filter.Tendsto
      (fun n : ℕ => (exp ((n : 𝕂)⁻¹ • A) *
                      exp ((n : 𝕂)⁻¹ • B)) ^ n)
      atTop (nhds (exp (A + B)))
\end{leancode}

Note the scalar action $\bullet$: the step is $n^{-1}$ \emph{acting on} $A$
in the $\mathbb{K}$-module $\A$, not a ring product with a coerced $n^{-1} \in
\A$. The two coincide in a unital algebra but are distinct terms in Lean, and
the derivative and integral lemmas we rely on are all stated for $\bullet$.

\paragraph{The \texttt{include} pattern.}
Since \texttt{NormedSpace.exp} no longer takes a field parameter in recent
\Mathlib, the variable $\mathbb{K}$ does not appear in the type of theorems
involving \texttt{exp}---in \texttt{norm\_exp\_le} below, both sides mention
only $\A$ and $\mathbb{R}$. Lean~4 drops unused type-level variables, so we
must write \texttt{include $\mathbb{K}$ in} before each lemma whose proof (but
not statement) requires $\mathbb{K}$
(e.g., \leanref{ExpBounds.lean}{72}{norm\_exp\_le}, whose proof rewrites with
\texttt{exp\_tsum\_form ($\mathbb{K}$ := $\mathbb{K}$)}):

\begin{leancode}
include 𝕂 in
lemma norm_exp_le (a : 𝔸) : ‖exp a‖ ≤ Real.exp ‖a‖ := by ...
\end{leancode}

\subsection{Telescoping in Lean}
\label{sec:lean-telescoping}

The telescoping identity~\eqref{eq:telescoping} is proved by induction on $n$
using \texttt{Finset.sum\_range\_succ} to peel off the last term
(\leanref{Telescoping.lean}{32}{telescoping\_direct}):

\begin{leancode}
theorem telescoping_direct (X Y : R) (n : ℕ) :
    (∑ k ∈ Finset.range n,
      X ^ k * (X - Y) * Y ^ (n - 1 - k)) =
      X ^ n - Y ^ n := by
  induction n with
  | zero => simp
  | succ m ih =>
    rw [Finset.sum_range_succ]
    -- Factor Y from the inner sum to get the
    -- inductive hypothesis
    ...
    noncomm_ring [pow_succ]
\end{leancode}

Here \texttt{R} is any \texttt{Ring}: the telescoping identity is purely
algebraic, so this file assumes no norm and no completeness.

The final step uses \texttt{noncomm\_ring}, which handles polynomial identities
in non-commutative rings. The standard \texttt{ring} tactic silently fails on
non-commutative goals---it produces an unsolved subgoal rather than an error,
which cost us significant debugging time early in the project.

\subsection{Exponential bounds: filling Mathlib gaps}
\label{sec:lean-exp-bounds}

The bound $\norm{e^a} \le e^{\norm{a}}$ was not available in \Mathlib for
general Banach algebras. Our proof works from the tsum representation
(\leanref{ExpBounds.lean}{72}{norm\_exp\_le}):

\begin{leancode}
include 𝕂 in
lemma norm_exp_le (a : 𝔸) : ‖exp a‖ ≤ Real.exp ‖a‖ := by
  rw [exp_tsum_form (𝕂 := 𝕂), real_exp_eq_tsum]
  calc ‖∑' n, (((Nat.factorial n : 𝕂)⁻¹) • a ^ n : 𝔸)‖
      ≤ ∑' n, ‖...‖         -- norm_tsum_le_tsum_norm
    _ ≤ ∑' n, ‖a‖ ^ n / n ! -- norm_exp_term_le
    _ = Real.exp ‖a‖        -- Real.hasSum_exp
\end{leancode}

The proof chains three standard ingredients: pulling the norm inside the tsum,
bounding each term $\norm{a^n/n!} \le \norm{a}^n/n!$ via submultiplicativity,
and recognizing the resulting series as $e^{\norm{a}}$.

\subsection{Step error: algebraic factorization in Lean}
\label{sec:lean-step-error}

The factorization~\eqref{eq:factorization} is verified by \texttt{noncomm\_ring}.
The cross-term bound~\eqref{eq:cross-term} is proved by induction using the
non-commutative recurrence
(\leanref{StepError.lean}{43}{norm\_pow\_add\_sub\_pow\_sub\_pow}):

\begin{leancode}
private lemma norm_pow_add_sub_pow_sub_pow (a b : 𝔸) :
    ∀ m : ℕ, 1 ≤ m →
      ‖(a + b) ^ m - a ^ m - b ^ m‖ ≤
        (‖a‖ + ‖b‖) ^ m - ‖a‖ ^ m - ‖b‖ ^ m := by
  intro m hm
  induction m with
  | zero => omega
  | succ n ih => ...
\end{leancode}

At 674~lines, \leansrc{StepError.lean} is the longest file in the basic
Lie--Trotter development (excluding commutator-scaling extensions).
The inductive cross-term bound accounts for roughly half of this length,
because each step requires explicit \texttt{norm\_mul\_le} applications and
\texttt{gcongr} for the inductive comparison.

\subsection{Commutator scaling: calculus in Lean}
\label{sec:lean-commutator}

The commutator-scaling proofs are the most technically demanding part of the
formalization, as they require doing analysis---derivatives, interval integrals,
FTC-2---in a non-commutative Banach algebra.

\paragraph{Derivative computations.}
\Mathlib's \texttt{hasDerivAt\_exp\_smul\_const'} supplies the
derivative $d/d\tau[e^{\tau B}] = B\,e^{\tau B}$, and the
product rule \texttt{HasDerivAt.mul} extends to normed algebras.
One subtlety: $(-\tau) \bullet B$ and $\tau \bullet (-B)$ are
definitionally distinct but equal; we must normalize before
applying derivative lemmas
(\leanref{CommutatorScaling.lean}{75}{hasDerivAt\_exp\_conj}):

\begin{leancode}
lemma hasDerivAt_exp_conj (A B : 𝔸) (s : ℝ) :
    HasDerivAt
      (fun u => exp (u • B) * A * exp ((-u) • B))
      (exp (s • B) * (B * A - A * B) *
       exp ((-s) • B)) s := by
  simp_rw [show ∀ u : ℝ, (-u) • B = u • (-B)
    from fun u => by rw [neg_smul, smul_neg]]
  ...
\end{leancode}

\paragraph{FTC-2 for integral representations.}
The conjugation identity~\eqref{eq:conj-ftc} and the integral error
representation~\eqref{eq:integral-error} are proved by
\texttt{integral\_eq\_sub\_of\_hasDerivAt}
(\leanref{CommutatorScaling.lean}{157}{exp\_conj\_sub\_eq\_integral}):

\begin{leancode}
theorem exp_conj_sub_eq_integral (A B : 𝔸) (τ : ℝ) :
    exp (τ • B) * A * exp ((-τ) • B) - A =
      ∫ s in (0 : ℝ)..τ,
        exp (s • B) * (B * A - A * B) *
        exp ((-s) • B) := by
  have hderiv : ∀ u ∈ Set.uIcc 0 τ,
      HasDerivAt ... := fun u _ =>
    hasDerivAt_exp_conj A B u
  have hint : IntervalIntegrable ... :=
    (continuous_exp_conj_deriv A B).intervalIntegrable 0 τ
  exact (integral_eq_sub_of_hasDerivAt hderiv hint).symm
\end{leancode}

\paragraph{Pulling constants through integrals.}
In a \texttt{NormedRing}, left-multiplication by a fixed element $c$ is
\emph{not} a scalar multiplication---it is ring multiplication. To derive
$c \cdot \int f = \int c \cdot f$, we use
\texttt{ContinuousLinearMap.intervalIntegral\_comp\_comm} with
$L = \texttt{ContinuousLinearMap.mul}\ \mathbb{R}\ \A\ c$, which views
left-multiplication as a continuous linear map.

\paragraph{Working over $\mathbb{R}$.}
The commutator-scaling files work over $\texttt{NormedAlgebra}\ \mathbb{R}\ \A$
directly, rather than a general $\mathbb{K}$. The reason is that the
\texttt{HasDerivAt}/\texttt{intervalIntegral} machinery requires
$\texttt{SMul}\ \mathbb{R}\ \A$, which is \emph{not} automatically synthesized
from $\texttt{[RCLike K] [NormedAlgebra K A]}$. For applications with
$\mathbb{K} = \mathbb{C}$, one can use
$\texttt{NormedAlgebra.restrictScalars}\ \mathbb{R}\ \mathbb{C}\ \A$.

\subsection{Multi-operator and Suzuki extensions}
\label{sec:lean-extensions}

The multi-operator formulas are proved by induction on the list of operators.
The palindromic product is defined recursively
(\leanref{MultiStrangCommutatorScaling.lean}{42}{palindromicProd}):

\begin{leancode}
def palindromicProd : List 𝔸 → 𝔸
  | [] => 1
  | [a] => exp a
  | a :: rest =>
    exp ((1/2 : ℝ) • a) *
    palindromicProd rest *
    exp ((1/2 : ℝ) • a)
\end{leancode}

The step size is not a parameter: it is folded into the operators by the
caller, which passes the list $[t \bullet A_1, \ldots, t \bullet A_m]$. This
keeps the recursion on a single argument, so the induction is on the list
alone. The inductive step peels off the outermost $e^{A_1/2}$ factors and
applies the two-operator Strang bound to the pair
$(A_1, A_2 + \cdots + A_m)$, then recurses on the inner product.

\section{The Fourth-Order Trotter Bound: Tighter Prefactors via BCH}
\label{sec:suzuki4}

The fourth-order Suzuki product formula $S_4(t)$ is the palindromic composition
of five Strang blocks with coefficients $(p, p, 1{-}4p, p, p)$, where
$p = 1/(4-4^{1/3})$ satisfies the Suzuki cubic $4p^3 + (1-4p)^3 = 0$:
\begin{equation}
  S_4(t)
    \;=\; S_2(pt)\,S_2(pt)\,S_2((1{-}4p)t)\,S_2(pt)\,S_2(pt),
  \qquad
  S_2(s) = e^{sA/2}\, e^{sB}\, e^{sA/2}.
\end{equation}
The Trotter error $\norm{S_4(t) - e^{t(A+B)}}$ vanishes to fourth order:
there exists a constant $C(A,B)$ such that
\begin{equation}
  \norm{S_4(t) - e^{t(A+B)}} \;\le\; C(A,B)\cdot t^5
\end{equation}
for all sufficiently small $t \ge 0$. The important question---answered
variously in the
literature with different prefactors---is: \emph{what is $C(A,B)$}?

Childs et al.~\cite{childsTheoryTrotterError2021} give, in the unitary
setting (skew-adjoint $A, B$, equivalently Hermitian generators evolving as
$e^{-itH}$) and for $t \ge 0$, an explicit bound of the
form
\begin{equation}\label{eq:childs-bound}
  \norm{S_4(t) - e^{t(A+B)}} \;\le\; t^5 \cdot \sum_{i=1}^{8} \alpha_i\, \norm{C_i(A,B)},
\end{equation}
where the $C_i$ are the eight four-fold nested commutators
$[X_1,[X_2,[X_3,[B,A]]]]$ with $X_j \in \{A, B\}$, and the coefficients $\alpha_i$
range from $0.0046$ to $0.0284$. This section discusses what we can (and cannot)
say about the tightness of these prefactors, how a Computer-Algebra-System
(CAS)-assisted BCH expansion yields strictly smaller coefficients, and how the
result is integrated with our Lean formalization.

\subsection{Childs et al.'s prefactors: rigorous, but not proven tight}
\label{sec:childs-heuristic}

A point of precision first: \cref{eq:childs-bound} is a \emph{rigorous}
upper bound. It is stated and proved as a proposition (Prop.~J.1 in the
arXiv version of~\cite{childsTheoryTrotterError2021}), derived via a
Duhamel expansion and the triangle inequality, and it holds for all
$t \ge 0$ in the unitary setting. What Childs et al.\ describe as
``heuristic'' are the \emph{strategies} used to make the prefactors small,
and what they leave unproven is tightness: ``we do not have a rigorous
proof of the tightness of these bounds.''

Concretely, they derive \cref{eq:childs-bound} by an elegant but
``balanced factoring,'' whose optimality is not proved, of the twelve-factor Duhamel
representation of the $S_4$ error integrand. The 4-fold commutator structure
emerges naturally from this Duhamel form; the specific numerical values of
$\alpha_i$ come from bounding the factored integrals term by term.

Two independent sources of looseness enter this derivation:

\paragraph{Over-completeness of the $C_i$ basis.}
The weight-five free Lie algebra on two generators has dimension six (Witt's
formula: $6 = \tfrac{1}{5}\sum_{d \mid 5} \mu(5/d)\,2^d$). Childs's eight
commutators $C_1,\ldots,C_8$ are therefore an \emph{over-complete} spanning
set: eight vectors in a six-dimensional space, satisfying two independent
linear relations. Whenever a Lie-algebra element such as the exact residual
$R_5$ is written as $\sum_i \beta_i C_i$, the signed
coefficients $\beta_i$ are not unique: one has two degrees of freedom
(corresponding to the two relations) to redistribute mass across the $C_i$'s.
Childs et al.'s $\alpha_i$ are \emph{not} coordinates of $R_5$ in this (or
any) representation---they bound the terms of the factored Duhamel integral
individually---so over-completeness enters their derivation only implicitly,
through which nested commutators the factoring happens to produce. For a
BCH-style bound it enters directly: each admissible representation yields a
bound $\sum_i \abs{\beta_i}\norm{C_i}$, and minimizing this weighted,
instance-dependent objective over representations (not an ordinary
coefficient $\ell^1$ norm) would generally give something smaller than any
fixed choice, including ours.

\paragraph{Balanced factoring vs.\ exact expansion.}
Childs's derivation bounds the Duhamel integral \emph{without} computing the
explicit Baker--Campbell--Hausdorff (BCH) expansion of $\log S_4(t)$. Each
term of the balanced factoring is bounded individually via the triangle
inequality. An exact BCH expansion gives a single expression for $R_5$ with
\emph{signed} coefficients whose cancellations have already been applied
(leveraging the Suzuki cubic condition and palindromic symmetry). Bounding
an expression with cancellations already applied is tighter than bounding
individual terms before cancellation.

Neither source of looseness affects validity---the $\alpha_i$ bound is
correct as stated---but both leave room below it, and that room is exactly
what the BCH expansion of \cref{sec:bch-expansion} occupies.

\subsection{The BCH expansion and its cancellations}
\label{sec:bch-expansion}

The BCH expansion of $\log S_4(t) - t(A+B)$ has only odd-order terms (by
palindromic symmetry), of which the leading one cancels under the Suzuki
condition:
\begin{equation}\label{eq:log-s4}
  \log S_4(t) - t(A+B) \;=\; \underbrace{0}_{O(t^2)} + \underbrace{0}_{O(t^3)}
  + \underbrace{0}_{O(t^4)} + t^5\,R_5(A,B) + O(t^7).
\end{equation}

\paragraph{$O(t^2)$ and $O(t^4)$ vanish by palindromic symmetry.}
The palindromic symmetry of $S_4$---invariance under time reversal
$\bar S_4(t) := S_4(-t)^{-1} = S_4(t)$---forces all even-order BCH
coefficients to vanish identically, regardless of $p$. This is Yoshida's
theorem for symmetric compositions.

\paragraph{$O(t^3)$ vanishes under the Suzuki cubic.}
The third-order BCH coefficient of $S_4$ has the form
$\sum_i c_i^3\, E_3(A,B)$, where $c_i \in \{p, p, 1{-}4p, p, p\}$ are the
block coefficients and
\begin{equation}\label{eq:E3}
  E_3 \;=\; \tfrac{1}{12}\comm{B}{\comm{B}{A}} \;-\; \tfrac{1}{24}\comm{A}{\comm{A}{B}}
\end{equation}
is the Strang BCH cubic coefficient. The Suzuki parameter $p = 1/(4-4^{1/3})$ is
precisely the root of
\begin{equation}
  \sum_{i=1}^{5} c_i^3 \;=\; 4p^3 + (1-4p)^3 \;=\; 0.
\end{equation}
Hence the $O(t^3)$ coefficient vanishes.

\paragraph{The $O(t^5)$ residual.}
The fifth-order residual $R_5(A,B)$ lives in the six-dimensional weight-5
Lie algebra. It is a \emph{specific} Lie polynomial, computable in closed
form from the BCH expansion of five palindromic exponentials.

\subsection{CAS computation and Lean integration}
\label{sec:cas}

Closed-form computation of $R_5$ requires tracking $\sim 30$ five-letter
monomials in $A$, $B$ through the BCH series---feasible but tedious by hand,
and natural for a Computer Algebra System.

We use a \textbf{symbolic non-commutative free algebra over $\mathbb{Q}[p]$},
implemented in Python with \texttt{sympy} (see
\leanref{../scripts/compute\_bch\_prefactors.py}{1}{scripts/compute\_bch\_prefactors.py}
in the repository):
\begin{enumerate}
  \item Elements of the free algebra are dictionaries mapping words
    (tuples over $\{A, B\}$) to rational-function coefficients in
    $\mathbb{Q}[p]$.
  \item $S_2(s)$ is computed by truncated power series for $e^{(s/2)A}$,
    $e^{sB}$, $e^{(s/2)A}$, composed non-commutatively.
  \item $S_4(t)$ is computed as the product of five $S_2$ blocks with
    coefficients $(p, p, 1{-}4p, p, p)$, truncated to total degree $5$.
  \item $\log S_4(t) - t(A+B)$ is computed by the truncated series
    $\log(1 + y) = y - y^2/2 + y^3/3 - y^4/4 + y^5/5$ with
    $y = S_4(t) - 1$.
  \item The Suzuki cubic reduction $p^3 = (48p^2 - 12p + 1)/60$
    is applied to every coefficient, reducing all powers of $p$ to
    polynomials of degree~$\le 2$.
  \item The script verifies degrees $2$, $3$, $4$ vanish identically and
    extracts the degree-5 part.
  \item The degree-5 residual is projected onto the Childs
    basis $\{C_1, \ldots, C_8\}$ via Gauss--Jordan elimination.
    Since the basis is over-complete, the projection has two free
    parameters; the script chooses the projection setting both to zero.
  \item Specializing at $p = 1/(4-4^{1/3}) \approx 0.4145$ gives numerical
    values for the eight signed coefficients $\beta_i(p)$; the Lean
    development stores the certified rational ceilings
    $\gamma_i \ge \abs{\beta_i(p)}$ at the $10^{-6}$ grid.
\end{enumerate}

\paragraph{The symbolic result.}
Before numerical specialization, the eight signed coefficients $\beta_i(p)$
are polynomials of degree~$\le 2$ in $p$:
\begin{equation}\label{eq:bch-gamma}
  \setlength{\arraycolsep}{1pt}
  \begin{array}{rlrlrl}
    \beta_1(p) &= \tfrac{127}{144000}\,p^2 + \tfrac{13}{36000}\,p - \tfrac{1}{24000}, \qquad &
    \beta_5(p) &= \tfrac{31}{9000}\,p^2 - \tfrac{13}{18000}\,p + \tfrac{1}{12000}, \\[3pt]
    \beta_2(p) &= \tfrac{1}{12000}\,p^2 + \tfrac{13}{6000}\,p - \tfrac{1}{4000}, &
    \beta_6(p) &= \tfrac{31}{3000}\,p^2 - \tfrac{13}{6000}\,p + \tfrac{1}{4000}, \\[3pt]
    \beta_3(p) &= 0, &
    \beta_7(p) &= 0, \\[3pt]
    \beta_4(p) &= -\tfrac{61}{9000}\,p^2 + \tfrac{13}{3000}\,p - \tfrac{1}{2000}, &
    \beta_8(p) &= \tfrac{1}{18000}\,p^2 + \tfrac{13}{9000}\,p - \tfrac{1}{6000}.
  \end{array}
\end{equation}
Note that $\beta_3$ and $\beta_7$ are identically zero---a consequence of
the particular projection onto the over-complete basis. Throughout,
$\beta_i(p)$ denotes the exact signed coefficients and $\gamma_i$ the
certified nonnegative rational ceilings $\gamma_i \ge \abs{\beta_i(p)}$
stored in the Lean development.

\paragraph{Comparison with Childs.}
Specializing at Suzuki $p$ and rounding up to the certified ceilings
$\gamma_i$, all non-zero $\gamma_i$ are strictly
smaller than the corresponding Childs coefficient $\alpha_i$, most by
an order of magnitude or more (\cref{tab:bch-vs-childs}).

\begin{table}[ht]
\centering
\begin{tabular}{@{}clrrl@{}}
\toprule
\textbf{Index} & \textbf{Commutator $C_i$} & \textbf{Ceiling~$\gamma_i$} & \textbf{Childs $\alpha_i$} & \textbf{Ratio} \\
\midrule
$C_1$ & $[A,[A,[A,[B,A]]]]$ & $0.000260$ & $0.0047$ & $18\times$ \\
$C_2$ & $[A,[A,[B,[B,A]]]]$ & $0.000663$ & $0.0057$ & $8.6\times$ \\
$C_3$ & $[A,[B,[A,[B,A]]]]$ & $0$        & $0.0046$ & $\infty$ \\
$C_4$ & $[A,[B,[B,[B,A]]]]$ & $0.000132$ & $0.0074$ & $56\times$ \\
$C_5$ & $[B,[A,[A,[B,A]]]]$ & $0.000376$ & $0.0097$ & $26\times$ \\
$C_6$ & $[B,[A,[B,[B,A]]]]$ & $0.001128$ & $0.0097$ & $8.6\times$ \\
$C_7$ & $[B,[B,[A,[B,A]]]]$ & $0$        & $0.0173$ & $\infty$ \\
$C_8$ & $[B,[B,[B,[B,A]]]]$ & $0.000442$ & $0.0284$ & $64\times$ \\
\bottomrule
\end{tabular}
\caption{Certified rational ceilings $\gamma_i \ge \abs{\beta_i(p)}$ vs.\ Childs et al.'s
coefficients $\alpha_i$ (at Suzuki $p = 1/(4-4^{1/3})$). The ``Ratio''
column is $\alpha_i/\gamma_i$---how many times tighter the certified
coefficient is. Entries $C_3$ and $C_7$ vanish exactly in our chosen
projection; other valid projections would redistribute their contribution
across the remaining six.}
\label{tab:bch-vs-childs}
\end{table}

\paragraph{Lean integration.}
The script emits a Lean snippet for the
\texttt{bchTightPrefactors} definition in
\leansrc{Suzuki4ViaBCH.lean}. The corresponding Lean theorem
\texttt{norm\_suzuki4\_level3\_explicit} states that there exist
$\delta > 0$ and $K' \ge 0$ such that, for all $0 \le t < \delta$,
\begin{equation}\label{eq:level3-sharp}
  \norm{S_4(t) - e^{t(A+B)}}
    \;\le\; t^5 \cdot \sum_{i=1}^{8} \gamma_i\, \norm{C_i}
    \;+\; K' \, t^6 .
\end{equation}
The leading coefficient is the CAS-certified $\gamma$-sum \emph{as it appears in
the statement}; only the $O(t^6)$ remainder is existentially quantified. The
theorem is fully proved at the project level: \texttt{\#print axioms
norm\_suzuki4\_level3\_explicit} returns only \texttt{[propext,
Classical.choice, Quot.sound]}, and it requires no C*-algebra or anti-Hermitian
structure---any complete normed algebra suffices.

The intermediate bridge theorem
\texttt{bch\_w4Deriv\_level3\_tight}---historically axiomatized as the
BCH-derived pointwise residual---is now a theorem composing the
five-factor palindromic quintic identification from the companion
\href{https://github.com/Jue-Xu/Lean-BCH}{Lean-BCH} project
(\texttt{suzuki5\_log\_product\_quintic\_tight\_at\_suzukiP}, revision
\texttt{05e8c52}) with the per-$i$ numerical bounds
$\abs{\beta_i(\text{suzukiP})} \le \gamma_i$ proved by \texttt{nlinarith}
on the tight rational interval $41449/100000 < \text{suzukiP} <
41450/100000$.

A companion theorem \texttt{norm\_suzuki4\_level3\_le\_childs\_pointwise}
formally verifies that the BCH-derived leading coefficient is dominated \emph{by}
Childs's---i.e.\ that our bound is the tighter of the two:
\begin{equation}
  t^5 \cdot \sum_{i} \gamma_i\,\norm{C_i}
    \;\le\;
  t^5 \cdot \sum_{i} \alpha_i\,\norm{C_i}.
\end{equation}
The element-wise comparison $\gamma_i \le \alpha_i$ is verified by
\texttt{nlinarith} from the explicit numerical values. Substituting it into
\cref{eq:level3-sharp} gives \texttt{norm\_suzuki4\_childs\_explicit}:
Childs et al.'s coefficients in the statement, still carrying the same
$K' t^6$ remainder---the Level~1 form described below. Sharpening further,
\texttt{norm\_suzuki4\_le\_childs\_near\_zero} shows that whenever the two
leading coefficients differ strictly, the BCH bound beats Childs et al.'s
\emph{without any remainder term} on a neighbourhood of the origin (the
radius $\delta$ is assembled from the existential Level-3 witnesses
$\delta_0, K'$, hence not explicit):
$\norm{S_4(t) - e^{t(A+B)}} \le t^5 \sum_i \alpha_i \norm{C_i}$ for
$0 \le t < \delta$.

\subsection{A local small-\texorpdfstring{$\tau$}{tau} refinement at order \texorpdfstring{$\tau^7$}{tau7}: the \texorpdfstring{$R_7$}{R7} correction}
\label{sec:r7}

A subtle question: is our leading-order Level 3 bound $t^5 \sum \gamma_i \norm{C_i}$
actually tighter than Childs's uniform bound, or does the comparison depend
on how we handle higher-order BCH contributions?

\paragraph{What the comparison does and does not establish.}
A natural worry is that Childs et al.'s larger coefficients $\alpha_i$ might
be \emph{uniformly tight} on some finite interval $[0, t^*]$, so that
$\alpha_i \cdot \norm{C_i}$ is the correct envelope over that interval
(folding in higher BCH orders), while our leading $\gamma_i$ is only an
asymptotic $t \to 0$ value. Two observations bear on this---neither of
which settles it:

\begin{enumerate}[leftmargin=*]
  \item \emph{The leading-order gap is large.} The ratio
    $\alpha_i / \gamma_i$ ranges from $8.6\times$ to $64\times$
    (\cref{tab:bch-vs-childs}). For the $\alpha_i$ to be uniformly tight on
    $[0, t^*]$, higher-order BCH contributions at $t = t^*$ would need
    to exactly saturate the gap $\alpha_i - \gamma_i$, and
    no mechanism in the balanced-factoring derivation---a
    triangle-inequality
    envelope of 12 individual Duhamel terms---targets this kind of
    saturation.

  \item \emph{Tightness is unproven, not disproven.} The original paper's
    statement ``we do not have a rigorous proof of the tightness of these
    bounds'' leaves uniform tightness \emph{open}, in both directions.
\end{enumerate}

What our development proves is confined to the leading coefficient
(termwise $\gamma_i \le \alpha_i$, so the $\gamma$-sum never exceeds the
$\alpha$-sum) and, under the strict-gap hypothesis, to a neighbourhood of
$0$ on which the full BCH bound undercuts the $\alpha$-envelope
(\texttt{norm\_suzuki4\_le\_childs\_near\_zero}). Global or uniform
optimality---of either family of coefficients---remains unknown.

\paragraph{A uniform BCH bound via \texorpdfstring{$R_7$}{R7}.}
To produce a rigorously uniform bound (not merely asymptotic), we extend
the CAS computation to order~7 (see
\leanref{../scripts/compute\_bch\_r7.py}{1}{compute\_bch\_r7.py}). The
script confirms:
\begin{itemize}[nosep]
  \item Orders $2$, $3$, $4$, and $6$ all vanish identically under
    palindromic symmetry + Suzuki cubic (extending the cancellations
    verified at order 5).
  \item Order $7$ has a non-zero residual $R_7(A, B)$ which, at Suzuki $p$,
    expands to 126 non-zero 7-letter words.
\end{itemize}
A straightforward triangle-inequality bound over these 126 words, using
$\norm{w} \le \max(\norm{A}, \norm{B})^7$ for any 7-letter word $w$, yields
the explicit uniform coefficient
\begin{equation}\label{eq:r7-K}
  K := \sum_{w} \abs{\mathrm{coef}(w)} \;\approx\; 0.01951,
  \qquad
  \norm{R_7(A, B)} \le K \cdot \max(\norm{A}, \norm{B})^7.
\end{equation}
The per-AB-degree breakdown of $K$ is shown in \cref{tab:r7-breakdown}.

\begin{table}[ht]
\centering
\begin{tabular}{@{}lrlr@{}}
\toprule
\textbf{AB signature} & $K_k$ \\
\midrule
$\norm{A}^1 \norm{B}^6$ & $0.000828$ \\
$\norm{A}^2 \norm{B}^5$ & $0.003835$ \\
$\norm{A}^3 \norm{B}^4$ & $0.006134$ \\
$\norm{A}^4 \norm{B}^3$ & $0.005510$ \\
$\norm{A}^5 \norm{B}^2$ & $0.002662$ \\
$\norm{A}^6 \norm{B}^1$ & $0.000540$ \\
\midrule
\textbf{Total $K$} & $0.019509$ \\
\bottomrule
\end{tabular}
\caption{R$_7$ uniform bound $K = \sum_k K_k$ split by AB signature. The
symmetric profile ($K_3, K_4$ peak) reflects palindromic structure.}
\label{tab:r7-breakdown}
\end{table}

\paragraph{The Level 4 existential-$\delta$ bound.}
Combining the certified ceilings $\gamma_i$ of the $R_5$ coefficients with
the $R_7$ uniform constant $K$, we obtain a bound on a small interval
$[0, \delta)$. The interval is a limitation of the proof route, not of the
statement: the bound is extracted from the BCH identification of
$\log S_4(t)$, which is available only within the convergence radius of the
BCH series, and $\delta$ records that radius. Duhamel-type analyses reach
all $t \ge 0$ directly---Childs et al.'s Prop.~J.1 is global in $t$ in the
unitary setting, and our compounded total-error bounds
(\cref{sec:s4-convergence}) hold, for every fixed $t$, for all sufficiently
large $n$ in a general normed algebra, at the price of an exponential
growth factor.
\begin{equation}\label{eq:level4}
\begin{aligned}
  &\exists\, \delta > 0,\ \exists\, C \ge 0,\ \forall\, t \in [0, \delta), \\
  &\quad \norm{S_4(t) - e^{t(A+B)}}
    \;\le\;
  C \cdot \Big( t^5 \cdot \sum_{i=1}^{8} \gamma_i \norm{C_i}
  \;+\;
  t^7 \cdot K \cdot \max(\norm{A}, \norm{B})^7
  \;+\;
  t^8 \Big).
\end{aligned}
\end{equation}
The multiplier $C$ absorbs the exp-Lipschitz inflation
$\exp(\norm{tH} + \norm{\delta_{\mathrm{bch}}})$ from lifting the BCH
log identification to $S_4$; the $t^8$ tail absorbs the higher-order
remainder of the BCH expansion ($R_9, R_{11}, \ldots$). For the small-$t$
regime of Trotter splitting these are negligible; if desired the same
pipeline can extend the CAS to arbitrary order with correspondingly
smaller corrections.

The Lean formalization adds the uniform constant
\leanref{Suzuki4ViaBCH.lean}{1116}{bchR7UniformConstant}
and the bound
\begin{equation}
  \norm{R_7} \;\le\; K \cdot \max(\norm{A}, \norm{B})^7,
\end{equation}
encoded as \leanref{Suzuki4ViaBCH.lean}{1122}{bchR7Bound}.
The Lean-BCH bridge corollary
\texttt{suzuki5\_log\_product\_septic\_at\_suzukiP} (now a theorem
upstream, with its $\tau^7$ Childs-7-basis identification factored
through two deeper Lean-BCH septic theorems at the pinned revision;
see \cref{sec:caveats}) composes with exp-Lipschitz on the Lean-Trotter
side to yield the theorem
\begin{leancode}
theorem norm_suzuki4_level4_uniform (A B : 𝔸)
    (hA : star A = -A) (hB : star B = -B) :
    let p : ℝ := 1 / (4 - (4 : ℝ) ^ ((1 : ℝ) / 3))
    ∃ δ > 0, ∃ C ≥ 0, ∀ τ : ℝ, 0 ≤ τ → τ < δ →
      ‖suzuki4Exp A B p τ - exp (τ • (A+B))‖ ≤
        C * (τ^5 * bchTightPrefactors.boundSum A B +
             τ^7 * bchR7Bound A B +
             τ^8)
\end{leancode}

\paragraph{Where the Childs comparison actually lives.}
It is tempting to compare \cref{eq:level4} against Childs's
$t^5 \sum \alpha_i \norm{C_i}$ directly, but \cref{eq:level4} is the
wrong place to do it: its multiplier $C$ is existentially quantified,
so the quantity $C \sum_i \gamma_i \norm{C_i}$ that sits in front of
$t^5$ is not a number the statement pins down. The comparison belongs
at Level~3, where \cref{eq:level3-sharp} certifies the leading
coefficient outright and only the $t^6$ remainder is existential.
Absorbing that remainder into the $\alpha$--$\gamma$ gap yields Childs's
\emph{own} bound with \emph{no} remainder term on a neighbourhood of the
origin, which is exactly
\leanref{Suzuki4ViaBCH.lean}{1387}{norm\_suzuki4\_le\_childs\_near\_zero}
(stated in \cref{sec:cas}): write $K' t^6 = t^5 \cdot (K' t)$ and observe
that $K' t$ falls below the gap $\sum_i (\alpha_i - \gamma_i) \norm{C_i}$
as soon as $t < \mathrm{gap}/(K'+1)$.

Its hypothesis, $\sum_i \gamma_i \norm{C_i} < \sum_i \alpha_i \norm{C_i}$,
compares two \emph{concrete} reals computed from the eight commutator norms,
so it is checkable for any given $A, B$; since $\gamma_i \le \alpha_i$
termwise, it fails only in the degenerate case where every commutator with
$\gamma_i < \alpha_i$ vanishes---e.g.\ $\comm{A}{B} = 0$, where both bounds
are $0$ and there is nothing to beat. And because it routes through Level~3
rather than Level~4, it needs neither anti-Hermiticity nor a C*-algebra.

\paragraph{What Level 4 is, and is not.}
L4 is not the bound one plugs a number into. Both $C$ and $\delta$ are
existentially quantified in its statement, so at a given step size $t$ one
knows neither whether $t < \delta$ nor what $C$ is; and since $C$ multiplies
the leading term as well, L4 certifies \emph{less} about the $t^5$ coefficient
than \cref{eq:level3-sharp} does. For comparing certified leading
coefficients---where the
$8.6$--$64\times$ reduction of the $\gamma_i$ error prefactors \emph{would}
translate into a fourth-root reduction in the required fourth-order steps
in the small-$t$ regime---the
statement to use is the sharp Level~3 bound \cref{eq:level3-sharp}, whose
leading coefficient is the CAS-certified $\gamma$-sum verbatim. Since
$\delta$ and $K'$ there (and $N$, $K'$ in the compounded bounds of
\cref{sec:s4-convergence}) remain existential, what is certified is an
explicit asymptotic leading coefficient, not an end-to-end numerical step or
gate count; extracting the latter would require re-proving those statements
with explicit constants.

What L4 does contribute is evidence that the CAS-assisted BCH pipeline
\emph{iterates}. It carries the order-$\tau^7$ residual through the full
chain---Lean-BCH septic identification, Childs-7-basis projection,
exp-Lipschitz lift---and lands the explicit constant
$K = 0.01951$ of \cref{eq:r7-K} inside a formal statement, with the next
correction pushed out to $\tau^8$. An analogous $R_9$ extension is mechanical.
The existential $\delta$ is a limitation of the logarithmic proof route
rather than of the statement: it records the convergence radius within which
the BCH identification of $\log S_4(t)$ is available, and it is not forced
by the conclusion---Childs et al.'s Duhamel-based Prop.~J.1 holds for all
$t \ge 0$ in the unitary setting, and the compounded bounds of
\cref{sec:s4-convergence} hold in a general normed algebra with an
exponential growth factor. (A polynomial right-hand side valid for all
$t \ge 0$ cannot be expected in a general normed algebra---the error may
itself grow exponentially in $t$, though it need not: it vanishes
identically for commuting $A, B$ and stays bounded in the unitary
setting---so global bounds there carry exponential factors.) The
existential $C$, by contrast, is an artifact
of how the exp-Lipschitz inflation was packaged, and is precisely what
\cref{eq:level3-sharp} avoids.

\subsection{Summary of the four BCH-based bounds}
\label{sec:summary-levels}

\cref{tab:levels} summarizes the four formalized bounds.

\begin{table}[ht]
\centering
\footnotesize
\setlength{\tabcolsep}{4pt}
\begin{tabular}{@{}l@{\hspace{4pt}}l@{\hspace{4pt}}l@{\hspace{4pt}}p{4.2cm}@{}}
\toprule
\textbf{Level} & \textbf{Lean theorem} & \textbf{Form} & \textbf{Notes} \\
\midrule
\textbf{L1 (Childs)}
  & \texttt{norm\_suzuki4\_le\_childs\_near\_zero}
  & $t^5 \sum \alpha_i \norm{C_i}$
  & Childs's own coefficients, with \emph{no} remainder term, on $[0,\delta)$.
    Corollary of L3 via $\gamma_i \le \alpha_i$; no axiomatization of
    the Childs bound.
    Needs the gap hypothesis
    $\sum \gamma_i\norm{C_i} < \sum \alpha_i\norm{C_i}$. \\
\midrule
\textbf{L2 (BCH unit)}
  & \texttt{norm\_suzuki4\_level2\_explicit}
  & $t^5 \sum \norm{C_i} + K' t^6$
  & Unit coefficients from primitive BCH
    $\abs{\beta_i(\text{Suzuki-}p)} \le 1$. Holds for every $p$
    satisfying the cubic condition. \\
\midrule
\textbf{L3 (BCH leading)}
  & \texttt{norm\_suzuki4\_level3\_explicit}
  & $t^5 \sum \gamma_i \norm{C_i} + K' t^6$
  & Certified $\gamma_i$, smaller than Childs's coefficient-by-coefficient
    (8.6--64$\times$); leading-order only. \\
\midrule
\textbf{L4 (BCH septic)}
  & \texttt{norm\_suzuki4\_level4\_uniform}
  & $C\,(t^5\!\sum \gamma_i \norm{C_i} + t^7 K M^7 + t^8)$
  & Small-$\tau$ bound, $M = \max(\norm{A}, \norm{B})$. $C$ absorbs the
    exp-Lipschitz inflation; being existential, it makes the certified
    $t^5$ coefficient weaker than L3's. \\
\bottomrule
\end{tabular}
\caption{Four formalized BCH-based bounds for $\norm{S_4(t) - e^{t(A+B)}}$.
Each row names the theorem whose \emph{type}
carries the stated Form; all are $\exists \delta > 0$ statements on $[0,\delta)$,
with $K'$ (and, in L4, $C$) existentially quantified but bound \emph{before} $\tau$.
L1 is a \emph{corollary} of L3---no axiomatization of the Childs bound on
Trotter's side---while L2/L3/L4 are Lean theorems composing the corresponding
Lean-BCH bridge corollaries with exp-Lipschitz lifts. At Lean-BCH
revision \texttt{05e8c52} (2026-07-28), all four levels are axiom-free
(\texttt{\#print axioms} returns only the standard foundational axioms),
including L4's $\tau^7$ extension.
Hypotheses also differ: L1--L3 hold in any complete normed algebra with
$\norm{1} = 1$, whereas L4 requires anti-Hermitian $A, B$ in a
C\textsuperscript{*}-algebra (\cref{sec:s4-convergence}).
}
\label{tab:levels}
\end{table}

\subsection{From step error to convergence:
  \texorpdfstring{$O(1/n^4)$}{O(1/n4)}}
\label{sec:s4-convergence}

Levels~1--4 all bound the \emph{single-step} error
$\norm{S_4(t) - e^{t(A+B)}}$. What a simulation pays for, however, is the
\emph{total} error after $n$ steps of size $t/n$, and the two are separated by
the telescoping amplification of \cref{sec:bg-trotter}: a step error of order
$\tau^{k+1}$ compounds to a total error of order $1/n^{k}$. Formalizing that
passage completes the hierarchy---the first-order and Strang formulas already
carry total-error theorems (\cref{sec:assembly,sec:strang}), whereas $S_4$
carried only a step bound.

\paragraph{The bound.}
Let $p$ satisfy the Suzuki cubic condition $4p^3 + (1-4p)^3 = 0$, as the
standard $p = 1/(4 - 4^{1/3})$ does. Then for all $A, B$ in a complete normed
algebra and every evolution time $t$, there exist a constant $C > 0$ and a
threshold $N$ such that
\begin{equation}\label{eq:s4-quartic}
  \norm{\bigl(S_4(t/n)\bigr)^n - e^{t(A+B)}} \;\le\; \frac{C}{n^4}
  \qquad \text{for all } n \ge N,
\end{equation}
and consequently $\bigl(S_4(t/n)\bigr)^n \to e^{t(A+B)}$ as $n \to \infty$
(\leanref{Suzuki4Convergence.lean}{158}{suzuki4\_total\_error\_quartic},
\leanref{Suzuki4Convergence.lean}{257}{suzuki4\_convergence\_quartic}).
Both are axiom-free. At the standard $p$ the cubic hypothesis is discharged
internally, leaving a statement with no hypotheses beyond the ambient algebra
(\leanref{Suzuki4Convergence.lean}{326}{suzuki4\_convergence\_quartic\_suzukiP}).

\paragraph{Proof.}
Telescoping (\cref{sec:telescoping}) gives
$\norm{X^n - Y^n} \le n\,\norm{X-Y}\,\max(\norm{X},\norm{Y})^{n-1}$
with $X = S_4(t/n)$ and $Y = e^{(t/n)(A+B)}$, and two ingredients feed it.
The step error $\norm{X-Y} \le C_0\,\abs{t/n}^5$ is the single-step BCH bound
(\leanref{Suzuki4BchBound.lean}{386}{exists\_norm\_s4Func\_sub\_exp\_le\_t5}),
valid once $\abs{t/n} < \delta$---which is exactly what the threshold $N$
enforces. The damping factor is a \emph{constant}: writing $K$ for the sum of
the magnitudes of the eleven $S_4$ coefficients weighted by $\norm{A}$ and
$\norm{B}$, we prove $\norm{S_4(\tau)} \le e^{\abs{\tau} K}$
(\leanref{Suzuki4Convergence.lean}{106}{norm\_suzuki4Exp\_le}), whence
\[
  \max(\norm{X},\norm{Y})^{n-1} \;\le\; \bigl(e^{\abs{t/n}\,K}\bigr)^{n}
  \;=\; e^{\abs{t}\,K},
\]
independent of $n$. The exponential target composes exactly,
$\bigl(e^{(t/n)(A+B)}\bigr)^n = e^{t(A+B)}$, so no error enters from that side,
and the pieces combine to
$n \cdot C_0 \abs{t}^5 n^{-5} \cdot e^{\abs{t} K} = O(1/n^4)$.

\paragraph{Scope of the bound.}
In \cref{eq:s4-quartic} both $C$ and $N$ are existential: $N$ marks where the
step size enters the BCH regime $\abs{\tau} < \delta$, and $C$ is inherited
from the crude single-step bound, whose constant is expressed through
$\norm{A}$ and $\norm{B}$. That establishes the \emph{order} but not the
\emph{constant}. Two things distinguish it from the generic $O(1/n^2)$ rate for
$S_4$ (\leanref{Suzuki4.lean}{376}{suzuki4\_error\_rate\_sq}), which holds for
\emph{every} $n \ge 1$ and needs \emph{no} cubic condition: fourth order is
precisely what the hypothesis $4p^3 + (1-4p)^3 = 0$ buys, and it is bought only
beyond the threshold $N$. (Both constants are existentially quantified; the
$O(1/n^2)$ proof's witness, $19 s^3 e^{6s} + 1$ with $s = \norm{A} + \norm{B}$,
is visible in the source but not in the statement.)

\paragraph{Commutator scaling in the total error.}
Compounding the \emph{sharp} step bound \eqref{eq:level3-sharp} through the same
telescope---rather than the crude one---keeps the certified prefactors alive all
the way to the $n$-step statement. At the Suzuki parameter, for $t \ge 0$ there
are $K' \ge 0$ and $N$ such that, for all $n \ge N$,
\begin{equation}\label{eq:s4-quartic-comm}
  \norm{S_4(t/n)^n - e^{t(A+B)}}
    \;\le\;
  e^{tK} \cdot \Bigl(\sum_{i=1}^{8} \gamma_i \norm{C_i}\Bigr)\,
    \frac{t^5}{n^4}
  \;+\; \frac{K'}{n^5},
  \qquad K = \texttt{s4Rate} + \norm{A} + \norm{B},
\end{equation}
and the same holds with Childs's $\alpha_i$ in place of $\gamma_i$
(\leanref{Suzuki4TightConvergence.lean}{59}{suzuki4\_total\_error\_commutator\_scaling},
\texttt{suzuki4\_total\_error\_childs\_scaling}; both axiom-free, and again with
no star structure). This is the commutator-scaling statement in the sense of
Childs et al.: the coefficient of $1/n^4$ is a sum of \emph{nested-commutator
norms}, so it collapses as $A$ and $B$ approach commutativity---behaviour that
an $\exists C$ bound such as \cref{eq:s4-quartic} cannot express.

\paragraph{No star structure is needed.}
\Cref{eq:s4-quartic} and the Level~1--3 step bounds alike hold in \emph{any}
complete normed algebra with $\norm{1} = 1$: they reach $S_4$ through BCH and
an exp-Lipschitz estimate, and never invoke the anti-Hermitian isometry
$\norm{e^{sX}} = 1$. Their Lean statements assume only \texttt{NormedRing},
\texttt{NormedAlgebra} $\mathbb{R}$, \texttt{NormOneClass} and
\texttt{CompleteSpace}. Among the $S_4$ results, only the Level~4 uniform bound
genuinely requires the C\textsuperscript{*}/anti-Hermitian setting---it carries
the hypotheses $A^* = -A$ and $B^* = -B$ explicitly---as does the Duhamel track
of \cref{sec:commutator}, where the isometry is what makes the FTC argument
close. The tight 4th-order Trotter bound is therefore available for arbitrary
bounded generators, not only for the skew-adjoint generators of unitary
dynamics.

\paragraph{One operator, two spellings.}
$S_4$ enters the development twice: as five Strang steps with a built-in $1/n$
(\texttt{suzuki4Step}, the object carrying the generic $O(1/n^2)$ rate), and as
eleven exponentials in a free step size $\tau$ (\texttt{suzuki4Exp}, the object
the BCH chain bounds). They are the same operator
(\leanref{Suzuki4Convergence.lean}{371}{suzuki4Step\_eq\_suzuki4Exp}); the four
junction merges of adjacent same-operator exponentials are supplied by the
five-block factorization of \cref{sec:bch-expansion}. Consequently
\eqref{eq:s4-quartic} transfers to \texttt{suzuki4Step}
(\leanref{Suzuki4Convergence.lean}{381}{suzuki4Step\_total\_error\_quartic}:
$\norm{\texttt{suzuki4Step}(n)^n - e^{A+B}} \le C/n^4$, i.e.\ the $t = 1$
instance, since \texttt{suzuki4Step} has the $1/n$ built into its definition;
general $t$ follows by rescaling $A \mapsto tA$, $B \mapsto tB$). The
improvement from $O(1/n^2)$ to $O(1/n^4)$ is therefore recorded on one and the
same object rather than inferred across two definitions.

\subsection{Remaining caveats}
\label{sec:caveats}

Two methodological caveats remain after the Level 4 uniform bound:

\paragraph{Caveat 1: Over-complete basis, non-unique projection.}
The Childs basis of 8 commutators is over-complete (the weight-5 free Lie
algebra is only 6-dimensional). Our projection sets two free parameters to
zero, assigning $\beta_3 = \beta_7 = 0$ (hence $\gamma_3 = \gamma_7 = 0$).
All such projections represent
the \emph{same} $R_5$ element, but the triangle-inequality upper bound
$\sum \abs{\beta_i} \norm{C_i}$ varies. The representation minimizing this
weighted, instance-dependent objective
gives the tightest bound in the Childs basis; finding it requires a linear
program. Our choice is \emph{a} valid projection, not provably \emph{the}
tightest in the Childs basis. It does happen to be strictly smaller than
Childs's on all eight coordinates---the Lean development verifies
$\gamma_i \le \alpha_i$ termwise for the certified ceilings---but that is a property of
\emph{this} projection, not of valid projections in general: adding a large
multiple of a relation vector yields another valid projection with
arbitrarily large coefficients.

\paragraph{Caveat 2: Dependency boundary and generated constants.}
Lean-Trotter carries zero own theorem-level BCH-interface axioms. All
four levels are \emph{fully proved at the project level}:
\texttt{\#print axioms norm\_suzuki4\_childs\_form\_via\_level3},
\texttt{\#print axioms norm\_suzuki4\_level2\_bch}, and \texttt{\#print
axioms norm\_suzuki4\_level3\_bch}, together with
\texttt{\#print axioms norm\_suzuki4\_level4\_uniform}, each return only
\texttt{[propext, Classical.choice, Quot.sound]} at Lean-BCH
revision~\texttt{05e8c52} (2026-07-28). The earlier revision
\texttt{d455ff0} discharged the previous B1.c quintic axiom; the current
pin also discharges the two former septic stepping stones.
The L2, L3 theorems compose, on the Lean-Trotter side, the Lean-BCH
bridge corollaries
\texttt{suzuki5\_log\_product\_quintic\_of\_IsSuzukiCubic} and
\texttt{suzuki5\_log\_product\_quintic\_tight\_at\_suzukiP}, each via
an exp-Lipschitz lift; L1 is a corollary of L3 via the Lean-proved
termwise inequality $\gamma_i \le \alpha_i$. The L4 small-$\tau$
($\tau^7$) refinement composes the bridge corollary
\texttt{suzuki5\_log\_product\_septic\_at\_suzukiP}; its symmetric-BCH
septic bridge and Childs-7-basis matching residual are both theorems in
the pinned companion
\href{https://github.com/Jue-Xu/Lean-BCH}{Lean-BCH} project.
The previously axiomatized
order-4 Taylor-coefficient identity
\texttt{bch\_iteratedDeriv\_s4Func\_order4}, which supports the abstract
$O(t^5)$ existence bound, is now a theorem via a three-slice chain
that composes a single-step $O(\abs{\tau}^5)$ bound from Lean-BCH
(\texttt{norm\_s4Func\_sub\_exp\_le\_of\_IsSuzukiCubic} composed with the
opaque-RHS corollary
\texttt{suzuki5\_bch\_M4b\_RHS\_le\_t5\_of\_IsSuzukiCubic}), a generic
Taylor-match-from-norm lemma
(\leanref{TaylorMatch.lean}{313}{iteratedDeriv\_eq\_of\_norm\_le\_pow}),
and our \texttt{iteratedDeriv\_exp\_smul\_mul\_at\_zero}. Axiom
inspection confirms that this theorem and the other headline results
(Lie--Trotter, Strang, single-step $S_4$ bound, all of L1/L2/L3)
depend only on \texttt{[propext, Classical.choice, Quot.sound]}, with
no transitive \texttt{sorryAx}. The four three-factor
symmetric-BCH-cubic identities previously axiomatized here (the
\texttt{symmetric\_bch\_cubic} family) are likewise theorems derived
from the corresponding Lean-BCH theorems specialized to
$\mathbb{K} = \mathbb{R}$. The $\gamma_i$ and $K$ literals are hard-coded
in the Lean sources and validated there only through the inequalities they
must satisfy (e.g.\ $\abs{\beta_i(\text{suzukiP})} \le \gamma_i$ via
\texttt{nlinarith} on tight rational brackets); the Python computation and
its transcription into Lean are \emph{not} themselves machine-checked.
The Lean theorem checks the inequalities actually used downstream, which is
the relevant trust boundary for the stated bounds.

\subsection{Summary}

This section demonstrates a \emph{CAS-assisted formalization workflow}
producing what is, to our knowledge, the first analytic tightening of the
Childs et al.\ (2021) $S_4$ prefactors (based on a literature search
documented in the project repository). Four bounds are formalized:

\begin{itemize}[leftmargin=*, nosep]
  \item \textbf{L1 (Childs form, axiom-free on Trotter's side).}
    Recovers, locally, the \emph{coefficient form} of Childs et al.'s
    (2021) Prop.~J.1: their coefficients $0.0046$--$0.0284$ with no
    remainder term on some neighbourhood of zero, under the strict-gap
    hypothesis of \cref{tab:levels}. (Their proposition itself is global in
    $t \ge 0$ in the unitary setting; L1 holds in any complete normed
    algebra, but only near $0$.) Derived from L3 via the
    Lean-proved termwise inequality $\gamma_i \le \alpha_i$;
    no axiomatization of Childs et al.'s bound is required.
  \item \textbf{L2 (BCH unit coefficients).} The simplest rigorous
    BCH-derived bound, using unit coefficients on all 8 Childs commutators.
    Follows from the primitive BCH bound $\abs{\beta_i(\text{Suzuki-}p)} \le 1$
    and is dominated numerically by L3.
  \item \textbf{L3 (BCH leading-order).} Replaces Childs et al.'s
    coefficients $\alpha_i$ with certified rational ceilings $\gamma_i$ of
    the signed BCH leading-order coefficients. All 8 are strictly smaller; two ($\gamma_3,
    \gamma_7$) vanish identically, and on the remaining six the improvement
    ratio ranges from $8.6\times$ (for $C_2, C_6$) to $64\times$ (for $C_8$).
    Its sharp form \cref{eq:level3-sharp} carries the $\gamma$-sum in the
    statement; with $\delta, K'$ existential, it certifies the asymptotic
    leading coefficient rather than an end-to-end numerical step count.
  \item \textbf{L4 (BCH small-$\tau$ refinement).} Adds an explicit $R_7$
    correction ($K \approx 0.01951$, also CAS-derived), pushing the first
    correction to the $\tau^5$ term out to order $\tau^7$. Both the
    multiplier $C$ and the radius $\delta$ are existentially quantified, so
    L4 is a demonstration that the CAS pipeline iterates to higher BCH
    order---not a bound one extracts a number from; that role belongs to
    \cref{eq:level3-sharp}.
\end{itemize}

Both the Lean formalization and the CAS computation are reproducible from
the project repository. Levels L1--L4 are fully Lean-verified without
project-specific axioms at Lean-BCH revision \texttt{05e8c52} (see
\cref{sec:caveats}). The three-factor symmetric BCH interface on which
L3/L4 depend is imported from Lean-BCH as a git dependency and is no
longer axiomatized on the Trotter side.

\section{Discussion}
\label{sec:discussion}

\subsection{AI assistance in the development process}
\label{sec:ai}

We used LLM-based coding assistants (Claude Code, Anthropic) to accelerate
formalization, in the same general direction as recent work on agentic
autoformalization~\cite{ren2026merlean}.

Our workflow followed a \emph{sorry-driven development} pattern:
\begin{enumerate}
  \item State all theorem signatures with \texttt{sorry} placeholders.
  \item Verify that the signatures compose correctly (the main theorem
    type-checks with all intermediates as \texttt{sorry}).
  \item Fill \texttt{sorry}s bottom-up, from leaves of the dependency DAG
    to the root.
\end{enumerate}
This approach worked well with LLM agents: independent leaves
(e.g., the four exponential bounds B1--B4) could be assigned to parallel
agents, each given the exact signature to fill and the list of available
lemmas.

\textbf{Where agents excelled.} ``Fill this \texttt{sorry} given these
lemmas''---tasks with a clear specification and known
ingredients---compiled on the first attempt when delegated to agents.
Most of the exponential bounds and the assembly step were completed this way.

\textbf{Where agents struggled.} ``Find the right proof approach''---tasks
requiring mathematical insight, such as the Strang commutator-scaling argument
or the algebraic factorization for the step error---required human-driven
exploration. The agent's tendency to attempt minor variations of a failing
approach (rather than reconsidering the overall strategy) became a bottleneck
for these tasks.

\textbf{A cautionary pattern.} As noted by Zhao et al.~\cite{zhao2025chsh},
LLM agents sometimes quietly strengthen hypotheses to make a proof trivial.
The \texttt{sorry}-first workflow guards against this: since all signatures are
fixed before proof filling begins, adding hypotheses breaks the type-checking
of downstream theorems.

\subsection{Mathlib ecosystem}
\label{sec:mathlib}

Several gaps in \Mathlib were encountered and filled during this project:
\begin{itemize}
  \item $\norm{e^a} \le e^{\norm{a}}$ for general Banach algebras
    (\leanref{ExpBounds.lean}{72}{norm\_exp\_le}).
  \item $e^x - 1 - x \le x^2/2 \cdot e^x$ for $x \ge 0$
    (\leanref{ExpBounds.lean}{129}{exp\_sub\_one\_sub\_bound\_real}).
  \item $\norm{e^a - 1} \le e^{\norm{a}} - 1$
    (\leanref{ExpBounds.lean}{88}{norm\_exp\_sub\_one\_le}).
\end{itemize}
These are natural candidates for upstreaming to \Mathlib.

Working over a general field $\mathbb{K}$ (via \texttt{RCLike}) introduced
friction with the calculus infrastructure, which assumes $\mathbb{R}$-module
structure. Our pragmatic solution---working over $\mathbb{R}$ directly in the
commutator-scaling files---is effective but means users must apply
\texttt{restrictScalars} for $\mathbb{C}$ applications. A unified treatment
would require additional typeclass instances in \Mathlib.

\subsection{Formalization-driven proof design}
\label{sec:formalization-insights}

The formalization process influenced proof design in several ways:

\textbf{Algebraic factorization.}
The standard approach to the step error---expanding $e^a\,e^b$ to second order
and bounding cross-terms---requires tracking non-commutative multinomial
coefficients. In Lean, this involves dozens of \texttt{norm\_mul\_le} and
\texttt{gcongr} steps for each cross-term. The algebraic factorization
(\cref{sec:step-error}) splits the problem into two independent bounds,
each handled by series comparison, reducing the proof length by approximately
half. This factorization is arguably cleaner than the standard approach even
on paper.

\textbf{Duhamel without Gronwall.}
The commutator-scaling proof was designed with Lean feasibility in mind.
The standard approach uses Gronwall's inequality or ODE uniqueness, which
would require formalizing existence and uniqueness for linear ODEs in
Banach spaces---theory not yet available in \Mathlib. The Duhamel/FTC approach
uses only the product rule and FTC-2, both already available.

\textbf{\texttt{noncomm\_ring} vs.\ \texttt{ring}.}
The \texttt{ring} tactic in Lean silently fails on non-commutative goals:
it produces an unsolved subgoal rather than an error message. We discovered
this only after extended debugging. The \texttt{noncomm\_ring} tactic handles
free non-commutative ring identities correctly but cannot use commutativity
hypotheses (e.g., $[a, e^a] = 0$). Our workaround: use \texttt{set} to make
exponentials opaque, rewrite commutativity as equalities, then apply
\texttt{noncomm\_ring} on the resulting free-ring identity.

\subsection{Relation to quantum simulation}
\label{sec:quantum}

Our results are stated for general complete normed algebras, which include
the bounded operators $B(\mathcal{H})$ on a Hilbert space as a special case.
Specialization to finite-dimensional matrix algebras
$M_d(\mathbb{C})$ is straightforward via \Mathlib's
\texttt{Matrix.instNormedRing}.

The commutator-scaling bounds are physically motivated by Hamiltonian
simulation: for a lattice Hamiltonian $H = \sum_j h_j$ where each $h_j$ acts
locally, the commutator $\comm{h_j}{h_k}$ vanishes for distant terms.
Our formalized bounds could serve as certified building blocks for verified
resource estimation in quantum algorithms, once combined with gate-count
analysis for individual exponentials.

\section{Conclusion and Future Work}
\label{sec:conclusion}

We have presented, to our knowledge, the first kernel-checked proof, with no
project-specific axioms, of norm convergence of the Lie--Trotter sequence for
arbitrary elements of a complete normed algebra, together with higher-order
extensions and commutator-scaling error bounds, in Lean~4 with the \Mathlib
library.
The development comprises over 14\,000 lines with zero \texttt{sorry} axioms
and zero own theorem-level BCH-interface axioms. Levels~1--4 are fully
proved at the project level: at Lean-BCH revision \texttt{05e8c52},
\texttt{\#print axioms} returns only \texttt{[propext,
Classical.choice, Quot.sound]} on both the $\tau^5$ headlines and the
Level~4 local small-$\tau$ ($\tau^7$) refinement. The development covers:
\begin{itemize}
  \item First-order Lie--Trotter convergence at rate $O(1/n)$;
  \item Symmetric Strang splitting at rate $O(1/n^2)$;
  \item Multi-operator generalizations of both formulas;
  \item First- and second-order commutator-scaling bounds via a
    Duhamel/FTC technique;
  \item A second-order Strang bound via norm-of-difference that captures
    signed cancellations invisible to the sum-of-norms form, and whose
    leading coefficient is machine-checked never to exceed it---strictly
    smaller whenever the two double commutators partially cancel, a gain
    measured at $9$--$50\%$ across standard spin-chain benchmarks
    (\cref{apd:tighter});
  \item A fourth-order Suzuki $S_4$ track with four CAS-assisted
    BCH-based bounds, including a replacement of the 2021 prefactors of
    Childs et al.\ (rigorous, but not proven tight) by CAS-certified
    coefficients smaller term by term by 8.6$\times$--64$\times$---the
    comparison machine-checked at the leading coefficient, with full-bound
    domination on a neighbourhood of zero under a strict-gap
    hypothesis---and an existential-$\delta$ small-$\tau$ refinement via
    an explicit $\tau^7$ correction;
  \item Fourth-order Suzuki convergence at rate $O(1/n^4)$, obtained by
    compounding the single-step $O(t^5)$ bound over $n$ steps --- and, from
    the sharp step bound, with an explicit commutator-scaled leading
    coefficient.
\end{itemize}
The first, second, and fourth-order convergence theorems together form a
complete verified hierarchy $O(1/n) \to O(1/n^2) \to O(1/n^4)$: for each
integrator the development supplies both a single-step error bound and the
compounded statement that the $n$-step product converges to $e^{t(A+B)}$ at the
corresponding rate.

The Duhamel approach to commutator scaling---avoiding ODE uniqueness in favor
of the fundamental theorem of calculus---may be of independent interest as a
proof strategy for related product formula bounds.

\paragraph{Future directions.}
Several extensions are natural:
\begin{enumerate}
  \item \emph{Mathlib contributions.}
    The exponential norm bounds (\leanref{ExpBounds.lean}{72}{norm\_exp\_le},
    \leanref{ExpBounds.lean}{88}{norm\_exp\_sub\_one\_le},
    \leanref{ExpBounds.lean}{129}{exp\_sub\_one\_sub\_bound\_real})
    are general-purpose lemmas suitable for upstreaming to \Mathlib.

  \item \emph{Trotter-native and higher-order BCH derivations.}
    The three-factor symmetric BCH identities, the two-factor
    palindromic quintic remainder, and the five-factor palindromic
    quintic and septic identifications are formalized in the Lean-BCH
    companion project at revision~\texttt{05e8c52}. A Trotter-native
    operator-algebra proof of the order-4 Taylor-coefficient identity
    (via \texttt{scripts/compute\_sumQuadCorr\_factored.py}) remains a
    useful independent derivation, while extending the certified BCH
    expansion to order nine would sharpen the local remainder.

  \item \emph{Specialization to quantum systems.}
    Connecting our abstract bounds to concrete gate-count estimates for
    Hamiltonian simulation on quantum computers would bridge the gap between
    formal mathematics and verified quantum algorithms.

  \item \emph{Unbounded operators.}
    The original Trotter--Kato theorem applies to strongly continuous
    semigroups on Banach spaces, covering unbounded generators.
    Formalizing this extension would require substantial additional
    functional analysis infrastructure in \Mathlib.
\end{enumerate}

More broadly, our project demonstrates that substantial results in operator
analysis---involving non-commutative algebra, power series manipulation,
Fr\'{e}chet calculus, and Bochner integration---are within reach of current
proof assistant technology.
We hope this work encourages further formalization of the mathematical
foundations underlying quantum algorithms.


\appendix
\section{Tighter Strang Bound via Norm-of-Difference}
\label{apd:tighter}

We show that the Strang commutator-scaling bound can be tightened by
delaying the triangle inequality.
The standard bound
(\cref{sec:strang}, following Childs et al.~\cite{childsTheoryTrotterError2021})
uses a \emph{sum of norms}:
\begin{equation}\label{eq:strang-sum}
  \norm{S_2(t) - e^{tH}}
    \le \frac{\norm{\comm{B}{\comm{B}{A}}}}{12}\,t^3
      + \frac{\norm{\comm{A}{\comm{A}{B}}}}{24}\,t^3.
\end{equation}
We prove a bound using the \emph{norm of a difference}, whose leading
coefficient is always $\le$ the sum-of-norms coefficient and strictly
smaller when the two double commutators partially cancel.
In the proportional model of \cref{eq:tighten-ratio} the cancellation
condition becomes positive alignment; in general no alignment is required,
and the numerical sweep of \cref{sec:comparison-standard} finds strictly
positive cancellation in every tested spin-chain configuration, including
at zero and negative alignment proxy.

\subsection{Setup}

Write $A' = A/2$, $H = A + B$, and recall the Strang splitting
$S_2(t) = e^{tA'}\,e^{tB}\,e^{tA'}$.
The Duhamel approach (\cref{sec:commutator}) gives
\begin{equation}\label{eq:strang-duhamel}
  S_2(t) - e^{tH} = e^{tH} \int_0^t \mathcal{T}_2(\tau)\,d\tau
  \qquad \text{(anti-Hermitian case: $\norm{e^{sX}} = 1$)},
\end{equation}
where $\mathcal{T}_2(\tau) = \text{Part}_A(\tau) + \text{Part}_B(\tau)$
is the Strang residual from the 4-factor product rule
(\cref{sec:strang}).

\subsection{Extracting the leading order}

The two parts have the following structure
(from the double FTC, \cref{sec:commutator}):

\paragraph{Part~A (from $A'$-conjugation).}
Define $g_1(s) = e^{sA'}\,\comm{A'}{B}\,e^{-sA'}$.
Then $\text{Part}_A(\tau) = \int_0^\tau (g_1(s) - g_1(\tau))\,ds$
(the ``subtract-constant-at-$\tau$'' trick).

Expanding $g_1(s) - g_1(\tau) = -\int_s^\tau g_1'(u)\,du$
where $g_1'(u) = e^{uA'}\,\comm{A'}{\comm{A'}{B}}\,e^{-uA'}$:
\begin{align}
  g_1(s) - g_1(\tau)
    &= -(\tau - s)\,\comm{A'}{\comm{A'}{B}} - \int_s^\tau
        \bigl(g_1'(u) - \comm{A'}{\comm{A'}{B}}\bigr)\,du
        \notag \\
    &= -(\tau - s)\,\comm{A'}{\comm{A'}{B}} + O(\tau^2).
        \label{eq:g1-expand}
\end{align}
The second term is bounded via the single FTC applied to $g_1'$:
\[
  g_1'(u) - \comm{A'}{\comm{A'}{B}}
    = \int_0^u e^{vA'}\,\comm{A'}{\comm{A'}{\comm{A'}{B}}}\,e^{-vA'}\,dv,
\]
so $\norm{g_1'(u) - \comm{A'}{\comm{A'}{B}}} \le
  \norm{\comm{A'}{\comm{A'}{\comm{A'}{B}}}}\,|u|$
(anti-Hermitian isometry).

Integrating \eqref{eq:g1-expand}:
\begin{equation}\label{eq:partA-expand}
  \text{Part}_A(\tau) = -\frac{\tau^2}{2}\,\comm{A'}{\comm{A'}{B}}
    + R_A(\tau),
\end{equation}
where $R_A(\tau)$ is a triple integral satisfying
\begin{equation}\label{eq:RA-bound}
  \norm{R_A(\tau)} \le
    \frac{\norm{\comm{A'}{\comm{A'}{\comm{A'}{B}}}}}{3}\,\tau^3.
\end{equation}
The bound \eqref{eq:RA-bound} follows from
$\norm{R_A(\tau)} \le \int_0^\tau \int_s^\tau |u|\,du\,ds
  \cdot \norm{\comm{A'}{\comm{A'}{\comm{A'}{B}}}}
  = \frac{\tau^3}{3}\,\norm{\comm{A'}{\comm{A'}{\comm{A'}{B}}}}$.

\paragraph{Part~B (from $B$-conjugation).}
The double-FTC remainder
$R_2(\tau) = e^{\tau B}\,A'\,e^{-\tau B} - A' - \tau\,\comm{B}{A'}$
satisfies (by the double-integral representation):
\begin{equation}\label{eq:R2-expand}
  R_2(\tau) = \frac{\tau^2}{2}\,\comm{B}{\comm{B}{A'}} + R_B(\tau),
  \qquad
  \norm{R_B(\tau)} \le
    \frac{\norm{\comm{B}{\comm{B}{\comm{B}{A'}}}}}{6}\,\tau^3.
\end{equation}
The bound on $R_B$ is exactly our triple-FTC result
(\texttt{norm\_exp\_conj\_sub\_comm2\_le\_of\_skewAdjoint}
in \leansrc{HigherCommutator.lean}).

Since $\text{Part}_B(\tau) = e^{\tau A'}\,R_2(\tau)\,e^{-\tau A'}$,
conjugation by $e^{\tau A'}$ introduces an additional correction:
\begin{equation}\label{eq:partB-expand}
  \text{Part}_B(\tau) = \frac{\tau^2}{2}\,\comm{B}{\comm{B}{A'}}
    + \underbrace{\frac{\tau^2}{2}\bigl(
      e^{\tau A'}\comm{B}{\comm{B}{A'}}e^{-\tau A'}
      - \comm{B}{\comm{B}{A'}}\bigr)}_{R_C(\tau)}
    + e^{\tau A'}\,R_B(\tau)\,e^{-\tau A'},
\end{equation}
where $R_C(\tau)$ is bounded by the single FTC on $\comm{B}{\comm{B}{A'}}$:
\begin{equation}\label{eq:RC-bound}
  \norm{R_C(\tau)} \le
    \frac{\norm{\comm{A'}{\comm{B}{\comm{B}{A'}}}}}{2}\,\tau^3.
\end{equation}

\subsection{The tighter bound}

Combining \eqref{eq:partA-expand} and \eqref{eq:partB-expand}:
\begin{equation}\label{eq:T2-refined}
  \mathcal{T}_2(\tau) = \frac{\tau^2}{2}\,D + R(\tau),
\end{equation}
where the \emph{effective double commutator} is
\begin{equation}\label{eq:D-def}
  \boxed{D \;=\; \comm{B}{\comm{B}{A'}} - \comm{A'}{\comm{A'}{B}}}
\end{equation}
and $R(\tau) = R_A(\tau) + R_C(\tau) + e^{\tau A'}\,R_B(\tau)\,e^{-\tau A'}$
satisfies
\begin{equation}\label{eq:R-bound}
  \norm{R(\tau)} \le T \cdot \tau^3,
  \qquad
  T = \frac{\norm{\comm{A'}{\comm{A'}{\comm{A'}{B}}}}}{3}
    + \frac{\norm{\comm{A'}{\comm{B}{\comm{B}{A'}}}}}{2}
    + \frac{\norm{\comm{B}{\comm{B}{\comm{B}{A'}}}}}{6}.
\end{equation}

Integrating:
\begin{align}
  \norm{S_2(t) - e^{tH}}
    &\le \int_0^t \norm{\mathcal{T}_2(\tau)}\,d\tau
    \le \int_0^t \Bigl(\frac{\norm{D}}{2}\,\tau^2 + T\,\tau^3\Bigr)\,d\tau
    \notag \\
    &= \frac{\norm{D}}{6}\,t^3 + \frac{T}{4}\,t^4.
    \label{eq:strang-tighter}
\end{align}

Substituting $A' = A/2$ (using $\comm{cX}{Y} = c\,\comm{X}{Y}$ for
central scalars, so $\comm{A'}{\comm{A'}{B}} = \tfrac{1}{4}\comm{A}{\comm{A}{B}}$
and $\comm{B}{\comm{B}{A'}} = \tfrac{1}{2}\comm{B}{\comm{B}{A}}$):

\begin{equation}\label{eq:strang-tighter-final}
  \boxed{
  \norm{S_2(t) - e^{tH}}
    \le \frac{1}{6}\left\|\frac{\comm{B}{\comm{B}{A}}}{2}
      - \frac{\comm{A}{\comm{A}{B}}}{4}\right\| t^3
    + \frac{T'}{4}\,t^4
  }
  \qquad (t \ge 0,\ A^* = -A,\ B^* = -B)
\end{equation}
where $T'$ is the triple-commutator sum \eqref{eq:R-bound} evaluated at
$A' = A/2$. As throughout this appendix, the bound is stated for skew-adjoint
$A, B$ in a C\textsuperscript{*}-algebra and for $t \ge 0$; these are exactly the
hypotheses of the Lean theorem.

\subsection{Comparison with the standard bound}
\label{sec:comparison-standard}

\paragraph{Relation.}
The standard bound \eqref{eq:strang-sum} can be written as
\[
  \frac{\norm{\comm{B}{\comm{B}{A}}}}{12}
  + \frac{\norm{\comm{A}{\comm{A}{B}}}}{24}
  = \frac{1}{6}\left(
    \frac{\norm{\comm{B}{\comm{B}{A}}}}{2}
    + \frac{\norm{\comm{A}{\comm{A}{B}}}}{4}
    \right).
\]
By the triangle inequality:
\[
  \left\|\frac{\comm{B}{\comm{B}{A}}}{2}
    - \frac{\comm{A}{\comm{A}{B}}}{4}\right\|
  \le
  \frac{\norm{\comm{B}{\comm{B}{A}}}}{2}
    + \frac{\norm{\comm{A}{\comm{A}{B}}}}{4},
\]
so the leading term of \eqref{eq:strang-tighter-final} is always
$\le$ the standard bound \eqref{eq:strang-sum}, and strictly smaller exactly
when the triangle inequality above is strict for the pair
$\bigl(\tfrac12\comm{B}{\comm{B}{A}},\, -\tfrac14\comm{A}{\comm{A}{B}}\bigr)$,
i.e.\ when the two double commutators partially cancel in $D$. How much
cancellation to expect is a property of the specific pair: the proportional
model at \cref{eq:tighten-ratio} below gives intuition, and the numerical
sweep at the end of this subsection measures it on standard spin chains.
This domination is itself
formalized: the displayed inequality is
\leanref{StrangCommutatorScalingTight.lean}{576}{norm\_D\_le\_sum\_of\_norms}
(axiom-free, no C*-structure), so the leading-coefficient domination is
machine-checked, not merely argued.

\paragraph{When is the tightening significant?}
The $O(t^4)$ correction in \eqref{eq:strang-tighter-final} is controlled
by triple commutator norms. Assuming $T' > 0$, for $t$ below the threshold
\[
  t^* = \frac{2}{3}\,
  \frac{\norm{\comm{B}{\comm{B}{A}}}/2 + \norm{\comm{A}{\comm{A}{B}}}/4
    - \norm{D}}{T'},
\]
the new bound \eqref{eq:strang-tighter-final} is tighter than
\eqref{eq:strang-sum}; if $T' = 0$ the correction term vanishes and the new
bound is at least as tight for every $t \ge 0$, and strictly tighter for
every $t > 0$ whenever
$\norm{D}$ lies strictly below the standard coefficient. In the Trotter
regime ($t = O(1/n)$ with
$n \to \infty$), the $O(t^4)$ correction is negligible and the tightening
from the leading term dominates.

\paragraph{How much tightening, and when.}
Write $X = \comm{B}{\comm{B}{A}}$ and $Y = \comm{A}{\comm{A}{B}}$, and suppose the
two double commutators are proportional, $Y \approx c\,X$ for a real $c$. Then
$D = \tfrac{X}{2} - \tfrac{Y}{4} = \bigl(\tfrac12 - \tfrac{c}{4}\bigr) X$, while the
standard bound's coefficient is $\tfrac{\norm{X}}{2} + \tfrac{\norm{Y}}{4}
= \bigl(\tfrac12 + \tfrac{\abs{c}}{4}\bigr)\norm{X}$. The tightening ratio is therefore
\begin{equation}\label{eq:tighten-ratio}
  r(c) \;=\;
  \frac{\norm{D}}{\norm{X}/2 + \norm{Y}/4}
  \;=\;
  \frac{\abs{2 - c}}{2 + \abs{c}} .
\end{equation}
Two consequences hold \emph{within this proportional model}. First,
$r(c) = 1$ for every $c \le 0$: proportional anti-alignment adds
constructively in the difference $D$, and the norm-of-difference bound then
gives no improvement. Second, the model's gain comes only from
\emph{positive} alignment, and it is total at $c = 2$ ---
i.e.\ $\comm{A}{\comm{A}{B}} \approx 2\,\comm{B}{\comm{B}{A}}$ makes $D = 0$ and
annihilates the entire $t^3$ term. A $37.5\%$ reduction ($r = 5/8$) occurs at
$c = 6/13 \approx 0.46$. The model is diagnostic, not a characterization:
generic pairs of double commutators are not scalar multiples of one another,
strictness of the triangle inequality does not require scalar alignment, and
the sweep below finds strictly positive gains even in instances whose
alignment proxy is zero or negative.

\paragraph{Numerical check on standard spin chains.}
A sweep over spin-$\tfrac12$ open chains---Heisenberg XXZ
($\Delta \in \{0.5, 1, 2\}$, even/odd bond split), transverse-field Ising
($h \in \{0.5, 1, 2\}$, both even/odd and ZZ/field splits), and
longitudinal-plus-transverse-field Ising (ZZ/field split), at
$L = 4, 6, 8, 10$ sites, spectral norms---finds a strictly positive gain in
\emph{every} one of the 40 configurations tested: $\norm{D}$ falls below the
standard coefficient
$\tfrac12\norm{\comm{B}{\comm{B}{A}}} + \tfrac14\norm{\comm{A}{\comm{A}{B}}}$
by $9$--$50\%$ (quoted for the better of the two operator orderings, both of
which are computed). Here $c_{\mathrm{eff}} =
\operatorname{Re}\langle X, Y\rangle_F / \langle X, X\rangle_F$, with
$X = \comm{B}{\comm{B}{A}}$ and $Y = \comm{A}{\comm{A}{B}}$, is the
least-squares coefficient of $Y$ on $X$ in the Frobenius inner
product---an unnormalized regression diagnostic, not a cosine; real pairs
are not scalar multiples of one another. Heisenberg chains
under the even/odd split show positive alignment
($c_{\mathrm{eff}} \approx +0.15$ to $+0.5$; gains
$39$--$50\%$). Across the tested configurations the gain is not
alignment-driven: in the
TFIM and LTFIM ZZ/field splits the two double commutators are exactly
Frobenius-orthogonal ($c_{\mathrm{eff}} = 0$, where the proportional model
predicts no gain), yet the measured gains are $24$--$33\%$, and
configurations with $c_{\mathrm{eff}} < 0$ retain $9$--$16\%$. Over the
tested sizes the gain decays slowly (Heisenberg XXX,
even/odd split: $50.0\%$, $42.7\%$, $41.3\%$, $38.6\%$ at $L = 4, 6, 8, 10$),
and the operator ordering matters (e.g.\ $15.5\%$ vs.\ $31.3\%$ for the TFIM
ZZ/field split at $h = 2$), so in practice one evaluates $\norm{D}$ for both
orderings and takes the better---cheap at benchmark sizes, though dense
spectral norms grow exponentially with $L$, so at scale one would instead
bound $\norm{D}$ through interaction-graph locality. Four representative
configurations were re-derived by an independent reimplementation (reversed
tensor-ordering convention, explicit-SVD norms), agreeing to relative error
below $10^{-9}$; scripts and raw
data accompany the repository
(\texttt{scripts/sweep\_strang\_alignment.py},
\texttt{claude/strang\_alignment\_sweep.csv}).

\subsection{Why this tightening is absent from Childs et al.}
\label{sec:why-absent}

The standard bound \eqref{eq:strang-sum} appears in
Childs et al.~\cite{childsTheoryTrotterError2021} (Prop.~16 in the arXiv
version, specialized to two summands) and
throughout the product formula literature. Several reasons explain why
the norm-of-difference formulation \eqref{eq:strang-tighter-final} has
not appeared previously.

\paragraph{Proof architecture.}
Childs et al.\ use a ``conjugation expansion'' (their Theorem~5) that
expands the error at each exponential interface independently, producing
separate terms for the $A$-$B$ and $B$-$A$ boundaries. The triangle
inequality on the resulting sum is the most direct bound within this
framework. The norm-of-difference requires an additional step:
recombining the interface terms, factoring out the common $\tau^2/2$
prefactor, and bounding the remainder with a \emph{triple}-commutator
estimate (\cref{eq:R-bound}). This extra step does not arise naturally
from the interface-by-interface analysis.

\paragraph{Framework compatibility.}
Childs et al.\ build a systematic framework around the
\emph{nested-commutator sum}
$\tilde\alpha_{\mathrm{comm}} =
  \sum \norm{\comm{H_{\gamma_{p+1}}}{\ldots\comm{H_{\gamma_2}}{H_{\gamma_1}}}}$,
which extends uniformly to any order~$p$ and any number of operators.
The effective commutator $D = \comm{B}{\comm{B}{A'}} - \comm{A'}{\comm{A'}{B}}$
is a \emph{difference} of two nested commutators---it does not fit into
the $\tilde\alpha_{\mathrm{comm}}$ framework as a single nested commutator.

\paragraph{Simplicity.}
The standard bound \eqref{eq:strang-sum} is a clean $t^3$ expression
with no correction term, yielding a simple Trotter-number formula
$r \sim (C\,t^3/\varepsilon)^{1/2}$ for resource estimation.
The tighter bound \eqref{eq:strang-tighter-final} introduces an $O(t^4)$
correction involving triple commutators, making the resource formula
more complex. For asymptotic scaling analysis, the simpler expression
is often preferred.

\paragraph{When the tightening matters.}
Despite these considerations, the leading coefficient of the
norm-of-difference bound is never larger, and the sweep of
\cref{sec:comparison-standard} realizes the improvement throughout the
tested spin-chain configurations: every one gains
$9$--$50\%$, including instances with zero or negative alignment proxy.
Given a specific
Hamiltonian $H = A + B$, the effective commutator $D$ is computable
directly at benchmark sizes (at scale, locality-based bounds replace the
dense computation), so one need not
guess: evaluate $\norm{D}$ for both operator orderings and compare. In the
Trotter regime ($t = O(1/n)$),
the $O(t^4)$ correction is subdominant, so the leading-term tightening
is what survives.
The formalization-driven insight here is that the Duhamel approach---by
producing the residual $\mathcal{T}_2(\tau)$ as a \emph{single}
algebraic expression---makes the norm-of-difference formulation natural,
whereas the interface-by-interface expansion does not.

\subsection{Extension to higher-order formulas}

The same principle extends to the fourth-order Suzuki formula $S_4$.
Childs et al.~\cite{childsTheoryTrotterError2021} (Prop.~J.1 in the arXiv
version) bound
the $S_4$ error as a sum of 8 separate 4-fold commutator norms:
\[
  \norm{S_4(t) - e^{tH}}
    \le \sum_{k=1}^{8} \alpha_k\,\norm{C_k}\,t^5,
\]
with numerical coefficients $\alpha_k \in [0.0046, 0.0284]$.
The norm-of-difference formulation would replace this with
\[
  \norm{S_4(t) - e^{tH}} \le \norm{R_5}\,t^5 + O(t^6),
  \qquad
  R_5 = \sum_{k=1}^{8} \beta_k(p)\,C_k,
\]
where $R_5$ is the single algebraic error expression of \cref{eq:log-s4} and
the $\beta_k(p)$ are its \emph{signed} BCH coefficients \eqref{eq:bch-gamma}.
A note on the tail: palindromic symmetry removes the $t^6$ term from the
\emph{logarithm} ($\log S_4(t) - tH = t^5 R_5 + O(t^7)$), but exponentiating
reintroduces one in the operator difference,
$S_4(t) - e^{tH} = t^5 R_5 + \tfrac{t^6}{2}(H R_5 + R_5 H) + O(t^7)$, so
$O(t^6)$ is what holds in a general normed algebra. For skew-adjoint $A, B$
the sharper $O(t^7)$ tail should be recoverable: $S_4(t)$ is then unitary,
$\Delta(t) = \log S_4(t) - tH$ is skew-adjoint for small $t$, and the
unitary Duhamel estimate $\norm{e^{tH+\Delta} - e^{tH}} \le \norm{\Delta}$
gives $\norm{S_4(t) - e^{tH}} \le \norm{R_5}\,t^5 + O(t^7)$; this local
skew-adjointness argument is not yet formalized or script-checked, so we
state the conservative tail.
Since $\norm{R_5} \le \sum_k \abs{\beta_k}\,\norm{C_k}
\le \sum_k \gamma_k\,\norm{C_k} \le \sum_k \alpha_k\,\norm{C_k}$
by the triangle inequality and \cref{sec:cas}, its leading $t^5$ coefficient
is never larger than that of the Level~3 bound or of the standard one, and
is strictly smaller whenever the 4-fold nested commutators $C_k$ partially
cancel.

A companion numerical test on the same spin-chain families as the sweep of
\cref{sec:comparison-standard} (60 rows; \texttt{scripts/test\_r5\_gain.py},
\texttt{claude/s4\_r5\_gain.csv}) measures precisely this gap:
$\norm{R_5}/\sum_k \gamma_k \norm{C_k}$ ranges from $0.14$ to $0.92$---a
further $8$--$86\%$ reduction beyond Level~3, with Heisenberg chains under
the even/odd split at $71$--$86\%$. Since all $\beta_k \ge 0$ at the Suzuki
point $p = 1/(4-4^{1/3})$, the gain is operator-level cancellation among
the $C_k$ up to sub-percent ceiling slack (the ratio against
$\sum_k \abs{\beta_k}\norm{C_k}$ differs by less than $0.001$). On the same instances the certified Level-3 sum
$\sum_k \gamma_k \norm{C_k}$ itself undercuts Childs et al.'s
$\sum_k \alpha_k \norm{C_k}$ by roughly $24$--$45\times$ (measured
$23.6$--$45.4\times$)---well beyond the worst-case
termwise guarantee---because the operators weight the coefficients
non-uniformly.

\subsection{Formalization status}

The triple-FTC building block (\cref{eq:R2-expand},
\leansrc{HigherCommutator.lean}) is fully formalized with 0
\texttt{sorry} axioms, and so is the tighter Strang bound
\eqref{eq:strang-tighter-final} itself:
\leanref{StrangCommutatorScalingTight.lean}{335}{norm\_strang\_comm\_scaling\_tight}
is sorry-free, and \texttt{\#print axioms} on it returns only
\texttt{[propext, Classical.choice, Quot.sound]}.
The \emph{comparison} with the standard bound is machine-checked as well:
\leanref{StrangCommutatorScalingTight.lean}{576}{norm\_D\_le\_sum\_of\_norms}
proves $\norm{D} \le \tfrac12\norm{\comm{B}{\comm{B}{A}}}
+ \tfrac14\norm{\comm{A}{\comm{A}{B}}}$ (equivalently
$\norm{D}/6 \le \norm{\comm{B}{\comm{B}{A}}}/12
+ \norm{\comm{A}{\comm{A}{B}}}/24$, \cref{sec:comparison-standard}), so the
tight bound's leading coefficient provably never exceeds the standard one; it
too is axiom-free and needs no C*-structure. What remains open is the
\emph{$S_4$ extension} of the norm-of-difference idea --- the $\norm{R_5}$-based
bound of \cref{sec:why-absent}: every formalized $S_4$ bound in this development
uses the sum $\sum_i \gamma_i \norm{C_i}$, never $\norm{R_5}$; the measured
$8$--$86\%$ gap reported in \cref{sec:why-absent} is what formalizing it
would recover.


\section*{Acknowledgments}
The Lean formalization was developed with assistance from Claude Code
(Anthropic) and Codex (OpenAI), used for proof search, repository review,
numerical verification, and manuscript preparation.

\bibliographystyle{plainnat}
\bibliography{references,ref_aps}
\end{document}
